% Noether Paper 43 complete zh-Hans-CN translated producer draft.
% Translation source: exact current-German standalone article, SHA-256
% 657799FA62D58538E6AFC810221DE2C9E1F7DC481E7DDEF2CAD76506DDEB8176.
% Inherited Simplified-Chinese interval is a drafting witness only, SHA-256
% 130646F67B105205CD783EDA2928A7FC45B14840D84D93DDC1AF9E1D725005CB.
% Status: translated producer draft; independent check pending.
% Compilation is mechanical production only. No source, semantic, formula,
% terminology, visual, regional, publication, or certification check is claimed.
\documentclass[11pt]{article}
\usepackage[a4paper,margin=2.4cm]{geometry}
\usepackage{fontspec}
\usepackage{xeCJK}
\usepackage{amsmath,amssymb,mathtools,array}
\setmainfont{Noto Serif}
\setCJKmainfont{Microsoft YaHei}
\setCJKsansfont{Microsoft YaHei}
\XeTeXlinebreaklocale "zh"
\XeTeXlinebreakskip=0pt plus 1pt
\setlength{\parindent}{2em}
\setlength{\parskip}{0.35em}
\setlength{\emergencystretch}{8em}
\allowdisplaybreaks
\raggedbottom
\makeatletter
\providecommand{\srcfn}[2]{%
  \begingroup
  \renewcommand{\thefootnote}{#1}%
  \footnote{#2}%
  \addtocounter{footnote}{-1}%
  \endgroup}
\providecommand{\srcfnmark}[1]{%
  \begingroup
  \renewcommand{\thefootnote}{#1}%
  \footnotemark%
  \addtocounter{footnote}{-1}%
  \endgroup}
\providecommand{\srcfntext}[2]{%
  \begingroup
  \renewcommand{\thefootnote}{#1}%
  \footnotetext{#2}%
  \endgroup}
\makeatother
\providecommand{\glqq}{“}
\providecommand{\grqq}{”}
\begin{document}
\setcounter{footnote}{0}

\section*{43. 理想微分与不同。}

\begin{center}
\emph{Journal f. d. reine u. angew. Math. 188 (1950), S. 1--21}\\[0.75ex]
\emph{Emmy Noether}（\(\dagger\)）著\srcfnmark{1)}。
\end{center}

\srcfntext{1)}{本文是 E. Noether 的遗稿，写于 1927/28 年冬，由 H. Grell 惠交编者。文中详尽阐述了作者 1929 年在布拉格德国数学家协会（D. M. V.）大会上所作的同名报告；见 Jahresbericht D. M. V. 39 (1930), S. 17 的斜体部分。从 § 6, 3 起，手稿显然原拟再作更细致的修订。不过，内行读者不难理解这里保留下来的提纲式论述；除少量润饰外，原稿未作改动。}

\subsection*{引言。}

众所周知，代数数域分歧理论的主定理断言：设一个素理想以第 \(\varrho\) 次幂出现在它所位于的素数 \(p\) 中，则不同（即分歧理想）至少能被该素理想的第 \((\varrho-1)\) 次幂整除；更确切地说，当 \(\varrho\) 不能被 \(p\) 整除时，恰能被第 \((\varrho-1)\) 次幂整除，而当 \(\varrho\) 能被 \(p\) 整除时，则能被更高次幂整除。

这个定理类似于如下事实：若一个线性因子以第 \(\varrho\) 次幂出现在多项式 \(f(x)\) 中，则 \(f(x)\) 的微分商 \(f'(x)\) 至少能被该线性因子的第 \((\varrho-1)\) 次幂整除；更确切地说，当 \(\varrho\) 不能被系数域的特征整除时，恰能被第 \((\varrho-1)\) 次幂整除，而当 \(\varrho\) 能被该特征整除时，则能被更高次幂整除。

下面我将说明，这里不只是形式上的类比。\emph{不同可以看作数域 \(K\) 的一个定义理想在一个相应的、以 \(x_1,\ldots,x_n\) 为不定元的整系数多项式环中的微分商；该微分商取在 \(x=\omega\) 处，也就是取 \(x_1=\omega_1,\ldots,x_n=\omega_n\)，其中 \(\omega_1,\ldots,\omega_n\) 表示 \(K\) 中全体整数所成系统---主阶 \(\mathfrak o\)---的一组模基。}这里，定义理想 \(\mathfrak M\) 对应于各 \(\omega\) 之间的全体关系，因而由所有满足 \(f(\omega)=0\) 的整系数多项式 \(f(x)\) 组成。微分商 \(\mathfrak M'[x\to\omega]\) 则定义为取在 \(x=\omega\) 处的差分商。

事实上，由于 \(f(\omega)=0\)，可将理想 \(\mathfrak M\) 看成由所有差 \(f(x)-f(\omega)\) 组成；再构造差分理想 \(\mathfrak B=(x_1-\omega_1,\ldots,x_n-\omega_n)\) 以及差分商 \(\mathfrak A=\mathfrak M:\mathfrak B\)。这里取商是指通常的理想商，并且是在以 \(K\) 中的整数、亦即 \(\mathfrak o\) 中的数为系数的 \(x\) 的多项式环内进行的。于是，不同 \(\mathfrak D\) 定义为
\[
\mathfrak D=\mathfrak A[x\to\omega].
\]
系数域扩张是必要的，否则差分理想和差分商根本无法定义。因此，这正是单不定元 \(x\) 的多项式微分商
\[
f'(\xi)=(f(x)-f(\xi)):(x-\xi),
\]
在 \(x=\xi\) 处取值的直接推广；这里也必须添入 \(\xi\) 来扩张多项式环，才能使各差有定义。事实上，只要各 \(\omega\) 的基由单个元素的各次幂组成，理想微分商也就过渡为多项式微分商。

这里给出的不同定义因而直接衔接于已知概念，但它引入了非不变的中间对象，因为理想 \(\mathfrak M\) 依赖于所选的基。若不用整系数多项式环，而引入一个与主阶 \(\mathfrak o\) 同构的环 \(\mathfrak O\)---它将与模 \(\mathfrak M\) 的剩余类环同构---并通过 \(\mathfrak o\) 把它的系数域扩张为 \(\mathfrak O_{\mathfrak o}\)，便可得到一个等价的、\emph{完全不变的定义}\srcfn{2)}{系数域扩张---即直积形成---在 § 1 中也针对非交换环加以处理，其一般性超过本文所需。对于本文的目的，§ 2 在交换环和有限基这一特殊化之下便已足够。}。这里，差分理想 \(\mathfrak B\) 由所有差 \((x-\xi)\) 导出，其中 \(x\) 和 \(\xi\) 在 \(\mathfrak O\) 与 \(\mathfrak o\) 的同构之下相互对应；差分商 \(\mathfrak A\) 等于零理想关于 \(\mathfrak B\) 的商；把 \(\mathfrak A\) 中的每个 \(x\) 都换成与它同构对应的 \(\xi\)，便得到不同。

下文将把这个不变定义置于首位，并由此直接定义任意交换环相对于一个子环的不同；这里只须满足关于系数扩张可能性的若干假设，例如在存在模基时这些假设总能满足，必要时可先过渡到某些商环。这样一来，特别是数域任意阶的不同便有了定义；一个阶模某个素数幂的剩余类环的不同也同样得到定义；此外还有相对不同，以及多不定元代数函数情形中的不同。不同的微分定义与通常定义的一致性，来自对数域通过直积形成而对应的伽罗瓦扩张环所作的一项精细结构考察；这项考察以直和分解为基础\srcfn{3)}{这里是对我关于判别式的论文所依据的思想的进一步发展，但不以读者熟悉该文为前提。参见 E. Noether, Der Diskriminantensatz für die Ordnungen eines algebraischen Zahl- und Funktionenkörpers, J. Math. 157 (1927), 82--104.}。在单不定元多项式的剩余类环这一特殊情形中，这项考察归结为对 Lagrange 插值公式的分析。更确切地说，我的意思如下：

引入一个与数域 \(K\) 同构并包含 \(\mathfrak O\) 的环 \(\mathfrak K\)，再用 \(K\) 的伽罗瓦域 \(F\) 扩张它的系数域---有理数。这个扩张后的系统 \(\mathfrak K_F\) 成为 \(n\) 个简单分量的直和；它们对应于 \(K\) 的 \(n\) 个同构（即向各共轭域的过渡），并且分别是差分商 \(\mathfrak A\) 及其各共轭的扩张。同时，\(\mathfrak K_F\) 的零理想成为由差分理想及其各共轭导出的 \(n\) 个理想的交。各简单分量具有 \(Fe^{(i)}\) 的形式，其中 \(e^{(i)}\) 表示单位元的一个分量；由于 \(e^{(i)}\) 在 \(x=\xi\) 时过渡为单位元，差分商等于 \(\mathfrak d e^{(i)}\)，这里 \(\mathfrak d\) 表示 \(\mathfrak o\) 的不同。不同 \(\mathfrak d\) 由体中所有满足如下条件的元素组成：该元素乘以 \(e^{(i)}\) 后属于 \(\mathfrak O_{\mathfrak o}\)，即该阶与自身的直积。对各共轭也有相应的结论。

为了由此得到通过互补模给出的定义，应注意 \(e^{(i)}\) 成为关于 \(\mathfrak O\) 或 \(\mathfrak K\) 中与各 \(\omega\) 相对应的基 \(x_1,\ldots,x_n\) 的线性形式；而各 \(e^{(i)}\) 中 \(x\) 的系数，恰好给出 \(\mathfrak o\)（及其各共轭）的互补模的一组基。由差分商 \(\mathfrak A\) 的系数属于 \(\mathfrak o\) 这一事实立即可知，不同等于 \(\mathfrak o\) 关于其互补模的商，从而得到 Dedekind 的定义。上述考察并不限于主阶；它们把每个阶的不同刻画为该阶关于其互补模的商。

在主阶的情形中，微分定义与通过基本方程给出的定义的一致性，来自上面提到的事实：一旦基由单个元素---这里即基本形式---的各次幂组成，理想微分商便过渡为多项式微分商。由此也就给出了后一种定义与 Dedekind 定义之间的概念联系。

在主阶的情形中，以基本方程为出发点，众所周知，可直接推出开头提到的主定理。为了精确确定指数，需要考察主阶模 \(p^t\) 的剩余类环的不同 \(\mathfrak d[p^t]\)，其中 \(t\) 足够大---取 \(t=n\) 即已足够，而对于二次域这也是必要的---；关键在于，把不同 \(\mathfrak d\) 模 \(p^t\) 考察时，它恰好过渡为 \(\mathfrak d[p^t]\)。这里仍可把考察限制于模 \(\mathfrak p^{t\varrho}\) 的剩余类环，因为环的直和对应于不同的直和。后一个环又与某个单不定元多项式的剩余类环同构，其系数模 \(p^t\) 取值；正如 Ö. Ore 已经证明的那样---他从另一途径得到这样一个多项式\srcfn{4)}{编者注：手稿中没有补全这一引文。参见本文仅具提纲性质的 § 7 中所列的提示。}---对这个多项式求微分，借助补充数便能精确统观一切可能的指数。

最后，这项结构考察还给出了如下著名定理：\(K\) 的一个上域的不同，等于相对不同与 \(K\) 的不同的乘积。这是因为同样的乘法性质对直和分解中的各单位元成立，因而也对各互补模成立。

通过过渡到某些商环，以上所述的全部事实对相对不同也完全成立，正如我在关于判别式的论文中为相对判别式所阐明的那样。再过渡到函数域之后，对于系数体特征为零的代数函数体，一切结论仍然成立；而在特征为 \(p\) 以及整代数函数的情形中，结构性质固然仍然保留（实质上甚至无须过渡到函数域），但分歧理论会因不可分扩张体的出现而发生变化，这里不再展开。我要指出，最重要的结果是上面概述的、对多不定元代数函数也成立的完全不变结构定理（§ 5 和 § 6）；过渡到互补模已经意味着把论证分解为系数恒等式，这一步只是为了把已知结果纳入其中。
% Source: Paper 43, printed pp. 4--6 (journal scan pp. 693--695)
\subsection*{§ 1. 环的系数域扩张或直积形成。定义理想\srcfnmark{5)}}
\srcfntext{5)}{对于他对本节提出的批评意见，我谨向 B. L. van der Waerden 致谢。}

本节以一般环为基础，并不预设乘法交换律。另一方面，为了简便起见，假定单位元存在；除非明确注明相反情形，这一假设在全文中均予采用。

\paragraph{1. 直积的定义。}
如果满足下列条件，环 \(\mathfrak O\times\mathfrak o\) 或 \(\mathfrak O_{\mathfrak o}\) 称为环 \(\mathfrak O\) 与 \(\mathfrak o\) 关于 \(\mathfrak h\) 的\emph{直积}，也称为\emph{由 \(\mathfrak O\) 通过把系数域 \(\mathfrak h\) 扩张为 \(\mathfrak o\) 而得到的环}\srcfn{6)}{下文大多采用第二种写法，用以提示系数域的扩张；第一种则是直积的通常名称。}：

\begin{enumerate}
\item[(1)] \(\mathfrak O_{\mathfrak o}\) 包含 \(\mathfrak O\) 和 \(\mathfrak o\) 作为子环，并且等于这两个环的乘积（也就是说，它等于 \(\mathfrak O_{\mathfrak o}\) 中所有同时包含 \(\mathfrak O\) 和 \(\mathfrak o\) 的环的交）。

\item[(2)] \(\mathfrak O\) 与 \(\mathfrak o\) 的交等于 \(\mathfrak h\)，其中 \(\mathfrak h\) 包含 \(\mathfrak O_{\mathfrak o}\) 的单位元。

\item[(3)] \(\mathfrak O\) 与 \(\mathfrak o\) 的元素两两可交换；因此，\(\mathfrak O_{\mathfrak o}\) 由全体双线性组合
\[
x_{i_1}y_{i_1}+\cdots+x_{i_r}y_{i_r}
\]
组成，其中各组 \(x\) 和 \(y\) 分别遍历 \(\mathfrak O\) 和 \(\mathfrak o\) 中任意有限多个元素所成的全部系统；而 \(\mathfrak h\) 位于 \(\mathfrak O\) 和 \(\mathfrak o\) 的中心，因而也位于 \(\mathfrak O_{\mathfrak o}\) 的中心。

\item[(4)] \(\mathfrak O_{\mathfrak o}\) 中的两个元素相等，当且仅当（当然，只要满足这一条件就一定相等）它们至少能以一种方式，借助一方面在 \(\mathfrak O\) 中、另一方面在 \(\mathfrak o\) 中成立的关系而形式地变为相等。其含义如下：若
\[
x_1y_1+\cdots+x_ny_n=0
\quad\hbox{在 }\mathfrak O_{\mathfrak o}\hbox{ 中},
\]
则可添入若干形式消失的组合 \((\gamma h)\delta-\gamma(h\delta)\) 之和，其中 \(h\) 属于 \(\mathfrak h\)，\(\gamma\) 属于 \(\mathfrak O\)，\(\delta\) 属于 \(\mathfrak o\)，从而把该式改写为
\[
(z_1\alpha_1+\cdots+z_r\alpha_r)+(t_1\beta_1+\cdots+t_s\beta_s)
\]
并满足 \(z_1=0,\ldots,z_r=0\)；\(\beta_1=0,\ldots,\beta_s=0\)。这里 \(\alpha,\beta\) 属于 \(\mathfrak o\)，\(z,t\) 属于 \(\mathfrak O\)，因而 \(z=0\) 表示 \(\mathfrak O\) 中各 \(x,\gamma,h\) 之间的关系，而 \(\beta=0\) 表示 \(\mathfrak o\) 中各 \(y,\delta,h\) 之间的关系。\srcfn{7)}{若不假定单位元存在，则 (3) 和 (4) 中还会出现附加的线性项。}
\end{enumerate}

\(\mathfrak O_{\mathfrak o}\) 中的两个元素相等，当且仅当它们的差以上述方式消失。

因此，由直积的要求，等式关系与运算关系（加法和乘法）都由各因子中的相应关系唯一确定。由此可知：若 \(\mathfrak O,\mathfrak o\) 及其交 \(\mathfrak h\) 分别同构于 \(\overline{\mathfrak O},\overline{\mathfrak o}\) 和 \(\overline{\mathfrak h}\)，并且直积 \(\mathfrak O\times\mathfrak o\) 存在，那么 \(\overline{\mathfrak O}\times\overline{\mathfrak o}\) 也存在，并且与 \(\mathfrak O\times\mathfrak o\) 同构。

设一个环 \(\mathfrak S\) 等于 \(\mathfrak O\) 与 \(\mathfrak o\) 的乘积，并且 \(\mathfrak O\) 与 \(\mathfrak o\) 的元素两两可交换，则 \(\mathfrak S\) 是 \(\mathfrak O\times\mathfrak o\) 的一个同态像，其中从 \(\mathfrak O\) 到 \(\overline{\mathfrak O}\) 以及从 \(\mathfrak o\) 到 \(\overline{\mathfrak o}\) 的映射仍为同构。因为，\(\mathfrak O\times\mathfrak o\) 中的等式蕴含在 \(\mathfrak S\) 中的等式之内，而在 \(\mathfrak O\) 或 \(\mathfrak o\) 上，它又分别与 \(\mathfrak O\) 或 \(\mathfrak o\) 在 \(\mathfrak O\times\mathfrak o\) 中的等式一致。\srcfn{8)}{另见 § 1, 4。}

\paragraph{2. 关于子环的直积存在的充要条件。}
设 \(\mathfrak O\) 与 \(\mathfrak o\) 是两个环，它们的交为位于二者中心的环 \(\mathfrak h\)。关于 \(\mathfrak h\) 的直积未必存在，下面的例子说明了这一点：令 \(\mathfrak h\) 为全体有理整数所成的环，令 \(\mathfrak O=\mathfrak h[x]\) 并满足 \(1-2x=0\)，再令 \(\mathfrak o=\mathfrak h[y]\) 并满足 \(1-2y=0\)，其中 \(y\) 表示一个新符号。这样，\(\mathfrak O\) 与 \(\mathfrak o\) 就是 \(\mathfrak h\) 的两个等价环，它们在集合论意义下的交为 \(\mathfrak h\)，而符号 \(x\) 与 \(y\) 之间没有给定任何运算关系。不可能存在一个同时包含 \(\mathfrak O\) 和 \(\mathfrak o\) 的环 \(R\)，使它们在 \(R\) 中---此时 \(x\) 与 \(y\) 之间已有运算关系---的交为 \(\mathfrak h\)。因为在 \(R\) 中有：
\[
0=(1-2x)y-x(1-2y)=y-x,
\]
而这里甚至没有用到 \(x\) 与 \(y\) 的可交换性。\srcfn{9)}{位于一个固定环内的交换环的商环---这里即添入 \(\frac12\)---本来就是唯一的。这类似于两个等价而相异的伽罗瓦域扩张不能嵌入同一个公共扩域。尽管如此，这里仍可通过直积形成而嵌入一个含零因子的环。}

\emph{直积存在的充要条件}如下：

\emph{至少必须存在一个同时包含 \(\mathfrak O\) 和 \(\mathfrak o\) 的环 \(R\)，使得 \(\mathfrak O\) 与 \(\mathfrak o\) 的交为 \(\mathfrak h\)，并使 \(\mathfrak O\) 与 \(\mathfrak o\) 的元素两两可交换。}

事实上，如果直积存在，它本身就是这样的一个环 \(R\)。反过来，为了构造直积，可以假定---下文始终作此假定---在 \(\mathfrak O\) 和 \(\mathfrak o\) 中不属于 \(\mathfrak h\) 的元素之间没有运算关系，因为通过过渡到 \(\mathfrak h\) 的等价扩张---即引入另一些符号---总能达到这一点。\srcfn{9a)}{参见 § 2, 1 及脚注 12)。}再设 \(\mathfrak h\) 包含 \(\mathfrak O\) 和 \(\mathfrak o\) 的单位元 \(e\)（不过也请参见注 7)）。现在把 \(\mathfrak T\) 定义为全体双线性组合
\[
x_{i_1}y_{i_1}+\cdots+x_{i_r}y_{i_r}
 =y_{i_1}x_{i_1}+\cdots+y_{i_r}x_{i_r}
\]
所成的集合，其中各组 \(x\) 和 \(y\) 分别遍历 \(\mathfrak O\) 与 \(\mathfrak o\) 中任意有限多个元素所成的全部系统，因而各 \(x\) 与各 \(y\) 在 \(\mathfrak T\) 中可交换。当且仅当两个这样的组合能借助 \(\mathfrak O\) 和 \(\mathfrak o\) 中的关系，按照 1, (4) 所精确定义的意义形式地变为相等时，才把它们规定为相等。再把
\[
\alpha_{i_1}x_{i_1}+\cdots+\alpha_{i_r}x_{i_r}
\quad\hbox{与}\quad
\beta_{i_1}x_{i_1}+\cdots+\beta_{i_r}x_{i_r}
\]
的和与积分别定义为
\[
(\alpha_{i_1}+\beta_{i_1})x_{i_1}+\cdots+(\alpha_{i_r}+\beta_{i_r})x_{i_r}
\quad\hbox{与}\quad
\sum_{\lambda,\mu}\alpha_{i_\lambda}\beta_{i_\mu}x_{i_\lambda}x_{i_\mu}.
\]
这里，在各 \(\alpha,\beta\) 之中也允许出现零。这个定义在上述等式意义下是唯一的，因为任何零关系（按 1, (4) 的意义）在右乘或左乘 \(\mathfrak o\) 或 \(\mathfrak O\) 中的一个元素之后仍然成为这样的零关系，而这些零关系的和与差也有同样性质。系统 \(\mathfrak T\) 构成一个同时包含 \(\mathfrak O\) 与 \(\mathfrak o\) 的环，并且在等式、加法和乘法方面满足直积的各项条件。作为一般嵌入假设的结果，\(\mathfrak T\) 也满足交的条件。事实上，令 \(\mathfrak S\) 为 \(R\) 中 \(\mathfrak O\) 与 \(\mathfrak o\) 的乘积；根据 1，\(\mathfrak S\) 是 \(\mathfrak T\) 的一个同态像，而在 \(\mathfrak O\) 与 \(\mathfrak o\) 上该同态为同构。若在 \(\mathfrak T\) 中有关系 \(x=y\)，其中 \(x\) 属于 \(\mathfrak O\)，\(y\) 属于 \(\mathfrak o\)，则在 \(\mathfrak S\) 中也有同一关系；因此，按假设，这两个元素在 \(\mathfrak S\) 中都属于 \(\mathfrak h\)，而由于映射在 \(\mathfrak O\) 和 \(\mathfrak o\) 上分别为同构，对 \(\mathfrak T\) 也有同样结论。因而，嵌入假设既必要又充分。

\paragraph{3. 环关于子环的定义理想。}
若 \(\mathfrak O\) 中的元素 \(\ldots,z,\ldots\) 所成的系统 \(S\) 满足：\(\mathfrak O\) 等于在 \(\mathfrak O\) 中由 \(\mathfrak h\) 与 \(S\) 生成的环---即等于 \(\mathfrak O\) 中所有同时包含 \(\mathfrak h\) 与 \(S\) 的环的交；用符号写作 \(\mathfrak O=\mathfrak h[S]\)---，则称 \(S\) 为 \(\mathfrak O\) 关于 \(\mathfrak h\) 的\emph{生成系}。仍假定 \(\mathfrak h\) 是 \(\mathfrak O\) 的中心的一个子环；于是，\(\mathfrak O\) 是非交换多项式环 \(\mathfrak h[\ldots Z\ldots]\) 的一个同态像，其中 \(Z\) 表示不定元，并且不定元系统与 \(S\) 具有相同的基数。这里所谓非交换多项式环，是指所有有限和 \(\sum h_iP_i=\sum P_i h_i\) 所成的系统，其中 \(P_i\) 是幂乘积，包括与指数 \(0\) 对应的单位元：
\[
P_i=Z_{i_1}^{\varrho_1}Z_{i_2}^{\varrho_2}\cdots Z_{i_n}^{\varrho_n}
\]
（当然，相同的指标可以重复出现），而等式定义为关于各个不同幂乘积的系数相等。由上述同态可知，\(\mathfrak O\) 与剩余类环 \(\mathfrak h[\ldots Z\ldots]/\mathfrak M\) 同构，其中双边理想 \(\mathfrak M\) 由所有满足如下条件的多项式 \(F(Z)\)（即 \(\mathfrak h[\ldots Z\ldots]\) 的元素）组成：\(F(z)\) 在 \(\mathfrak O\) 中消失；这里 \(Z\) 与 \(z\) 在该同态之下相互对应。\emph{理想 \(\mathfrak M\) 称为关于 \(\mathfrak h\) 的一个定义理想，更确切地说，是与生成系 \(z\) 相对应的理想。}由定义可知，若两个环关于同构对应的子环 \(\mathfrak h,\overline{\mathfrak h}\) 和生成系 \(S,\overline S\) 同构，那么它们的定义理想也同构；特别地，当 \(\overline{\mathfrak h}\) 与 \(\mathfrak h\) 重合时，还可把这两个定义理想视为相同。

因为已经假定 \(\mathfrak h\) 位于 \(\mathfrak O\) 的中心，所以只要 \(\mathfrak O\) 有一个仅由单个元素 \(z\) 组成的生成系 \(S\)，\(\mathfrak O\) 就必为交换环。若 \(\mathfrak M\) 还是主理想，则称由 \(\mathfrak M\) 的一个基多项式 \(G(Z)\) 给出的等式 \(G(z)=0\) 为 \(\mathfrak O\) 关于 \(\mathfrak h\) 的一个定义方程。

\paragraph{4. 直积的定义理想。}
设（如 3 中）\(\mathfrak O\) 与 \(\mathfrak h[\ldots Z\ldots]/\mathfrak M\) 同构，其中 \(\mathfrak M\) 是定义理想。在 \(\mathfrak h[\ldots Z\ldots]\) 之外，再考察多项式环 \(\mathfrak o[\ldots Z\ldots]\)，即全体有限和
\[
\sum \gamma_iP_i=\sum P_i\gamma_i
\qquad\hbox{其中 }\gamma_i\hbox{ 属于 }\mathfrak o
\]
所成的系统，并把等式定义为系数相等。\(\mathfrak o[Z]\) 包含 \(\mathfrak h[Z]\)，并且是 \(\mathfrak h[\ldots Z\ldots]\) 与 \(\mathfrak o\) 关于 \(\mathfrak h\) 的直积；换言之，它由 \(\mathfrak h[\ldots Z\ldots]\) 通过扩张系数域而得到。\srcfn{10)}{因为等式关系与运算关系都得到满足，并且由 \(\gamma=h+h_1P_1+\cdots+h_rP_r\)，按照 \(\mathfrak o[\ldots Z\ldots]\) 中的等式可知 \(\gamma=h\)。对此可参见 § 2, 2。本项并不使用 \(\mathfrak o[\ldots Z\ldots]\) 是直积这一事实。}以 \(\mathfrak M_{\mathfrak o}\) 表示 \(\mathfrak M\) 关于 \(\mathfrak o[\ldots Z\ldots]\) 的扩张理想---即 \(\mathfrak o[\ldots Z\ldots]\) 中所有包含 \(\mathfrak M\) 的双边理想的交---；因而，它是 \(\mathfrak M\) 中各元素以 \(\mathfrak o\) 中元素为系数的全体线性组合所成的系统，并且由于 \(\mathfrak o\) 与各 \(Z\) 可交换，它已经是一个双边理想。\emph{如果关于 \(\mathfrak h\) 的直积 \(\mathfrak O_{\mathfrak o}\) 存在，那么 \(\mathfrak O_{\mathfrak o}\) 与 \(\mathfrak o[\ldots Z\ldots]/\mathfrak M_{\mathfrak o}\) 同构。}事实上，\(\mathfrak O_{\mathfrak o}\) 是 \(\mathfrak o[\ldots Z\ldots]\) 的一个同态像，因而与模 \(\mathfrak O_{\mathfrak o}\) 关于 \(\mathfrak o\) 的定义理想所成的剩余类环同构，该定义理想对应于生成系 \(z\)。而根据 1, (4) 的等式定义，\(\mathfrak O_{\mathfrak o}\) 中各 \(z\) 之间的关系可以写成 \(\mathfrak O\) 中已有关系以 \(\mathfrak o\) 中元素为系数的线性组合；所以，\(\mathfrak M_{\mathfrak o}\) 就是定义理想。

\paragraph{注。}
同样的考察还给出直积及其相应定义理想的一种对称但较为繁复的表示。设 \(\mathfrak o\) 与 \(\mathfrak h[\ldots Y\ldots]/\mathfrak N\) 同构，其中 \(Y\) 表示不同于 \(Z\) 的符号，从而多项式环 \(\mathfrak h[\ldots Y\ldots Z\ldots]\) 有定义。于是，\(\mathfrak O_{\mathfrak o}\) 同构于
\[
\mathfrak h[\ldots Y\ldots Z\ldots]/(\mathfrak M,\mathfrak N,\ldots,YZ-ZY,\ldots),
\]
其中右侧的定义理想由 \(\mathfrak M,\mathfrak N\) 和所有组合 \(YZ-ZY\) 生成。（后一类组合表示 \(\mathfrak o\) 与 \(\mathfrak O\) 中元素的可交换性，前两者表示 \(\mathfrak O_{\mathfrak o}\) 中的关系都是 \(\mathfrak O\) 和 \(\mathfrak o\) 各自关系的推论。）
% Source: Paper 43, printed pp. 7--10 (journal scan pp. 696--699)
\subsection*{§ 2. 具有独立模基的环。直积的构造。}

\paragraph{1. 直积的构造。}
若 \(\mathfrak O\) 中的元素 \(\ldots,t,\ldots\) 所成的系统 \(T\) 满足：\(\mathfrak O\) 等于在 \(\mathfrak O\) 中由 \(T\) 导出的 \(\mathfrak h\)-模，则称 \(T\) 为 \(\mathfrak O\) 的一组 \(\mathfrak h\)-模基。由于单位元存在，\(\mathfrak O\) 此时由全体线性组合 \(h_{i_1}t_{i_1}+\cdots+h_{i_r}t_{i_r}\) 组成。若
\[
h_{i_1}t_{i_1}+\cdots+h_{i_r}t_{i_r}=0
\quad\hbox{总能推出}\quad
h_{i_1}=0,\ldots,h_{i_r}=0,
\]
则称这组模基是独立的。\emph{如果 \(\mathfrak O\) 有一组包含单位元的独立模基，那么，对于 \(\mathfrak h\) 的任意扩张环 \(\mathfrak o\)，只要 \(\mathfrak h\) 是其中心的一个子环，并且 \(\mathfrak O\) 与 \(\mathfrak o\) 的集合论意义下的交 \([\mathfrak O,\mathfrak o]\) 等于 \(\mathfrak h\)，直积 \(\mathfrak O_{\mathfrak o}\) 就存在。}\srcfn{11)}{与此对称，也可以对 \(\mathfrak o\) 作出基的假设。根据 § 2, 2，在不失一般性之下---参见 E. Noether, Der Diskriminantensatz für die Ordnungen eines algebraischen Zahl- oder Funktionenkörpers, dies. Journ. 157 (1927), 82--104, §§ 1, 2---假定 \(\mathfrak O\) 与 \(\mathfrak o\) 中不属于 \(\mathfrak h\) 的元素之间不存在运算关系。}根据 § 1, 2，只须构造一个相应的嵌入环 \(\mathfrak R\)；在这里，该环随后就会证明为直积本身。把 \(\mathfrak R\) 定义为全体线性组合
\[
\gamma_{i_1}t_{i_1}+\cdots+\gamma_{i_r}t_{i_r}
=t_{i_1}\gamma_{i_1}+\cdots+t_{i_r}\gamma_{i_r}
\]
所成的系统，其中各 \(t_{i_\nu}\) 遍历独立基 \(T\) 中任意有限多个元素所成的全部系统，相应的 \(\gamma\) 取自 \(\mathfrak o\)。等式定义为关于各 \(t_{i_\nu}\) 的系数相等；和与积则按通常的运算法则定义，如 § 1, 2 所述。由基的性质，任意两个 \(t_{i_\nu}\) 的积都是有限多个基元素以 \(\mathfrak h\) 中元素为系数的线性组合，所以 \(\mathfrak R\) 成为一个环。\(\mathfrak R\) 包含 \(\mathfrak O\)，因为 \(\mathfrak h\) 是 \(\mathfrak o\) 的子环；它也包含 \(\mathfrak o\)，因为 \(T\) 包含单位元。又因为各 \(\gamma\) 分别与各 \(h_{i_\nu}\) 以及各 \(t_{i_\nu}\) 两两可交换，\(\mathfrak O\) 与 \(\mathfrak o\) 在 \(\mathfrak R\) 中的元素两两可交换。最后，\(\mathfrak O\) 与 \(\mathfrak o\) 在 \(\mathfrak R\) 中的交为 \(\mathfrak h\)，因为关系
\[
\gamma=\gamma e=h_1t_1+\cdots+h_rt_r
\]
按照等式定义---并利用 \(e\) 按假设是基的一个元素---说明其中一个 \(t\) 等于 \(e\)，其余项消失，因而 \(\gamma=h\)。所以直积存在；它等于 \(\mathfrak R\)，因为 \(\mathfrak R\) 等于 \(\mathfrak O\) 与 \(\mathfrak o\) 的乘积，并满足可交换性与等式条件（由于基元素独立，这个条件归结为系数相等）。\(\mathfrak O\) 的这组基 \(T\) 同时成为 \(\mathfrak O_{\mathfrak o}\) 的一组 \(\mathfrak o\)-基。

\paragraph{2. 扩张模与收缩模。}
若 \(\mathfrak B\) 是 \(\mathfrak O\) 中的一个 \(\mathfrak h\)-模，则其扩张 \(\mathfrak B_{\mathfrak o}\) 是指 \(\mathfrak O_{\mathfrak o}\) 中的这样一个 \(\mathfrak o\)-模：它等于 \(\mathfrak O_{\mathfrak o}\) 中所有包含 \(\mathfrak B\) 的 \(\mathfrak o\)-模的交。若 \(\mathfrak C\) 是 \(\mathfrak O_{\mathfrak o}\) 中的一个 \(\mathfrak o\)-模，则它在 \(\mathfrak O\) 中的收缩模定义为交 \([\mathfrak C,\mathfrak O]\)。若 \(\mathfrak B\) 同时是 \(\mathfrak O\) 中的理想，则 \(\mathfrak B_{\mathfrak o}\) 是 \(\mathfrak O_{\mathfrak o}\) 中的理想；若 \(\mathfrak C\) 是 \(\mathfrak O_{\mathfrak o}\) 中的理想，则 \([\mathfrak C,\mathfrak O]\) 是 \(\mathfrak O\) 中的理想。

\paragraph{定理。}
\emph{如果 \(\mathfrak O\) 中的一个模 \(\mathfrak B\) 有一组独立 \(\mathfrak h\)-模基，而且这组基可以补成 \(\mathfrak O\) 的一组独立 \(\mathfrak h\)-模基，那么 \(\mathfrak B\) 是它在 \(\mathfrak O_{\mathfrak o}\) 中的扩张模 \(\mathfrak B_{\mathfrak o}\) 的收缩。}

\paragraph{证明。}
以 \(T\) 表示 \(\mathfrak B\) 的一组独立 \(\mathfrak h\)-模基；按假设，可用 \(\mathfrak O\) 中由元素 \(\overline t\) 组成的系统 \(\overline T\) 把它补成 \(\mathfrak O\) 的一组独立 \(\mathfrak h\)-模基。于是，\(T\) 同时是 \(\mathfrak B_{\mathfrak o}\) 的一组独立 \(\mathfrak o\)-基，而 \(T\) 与 \(\overline T\) 合在一起则是 \(\mathfrak O_{\mathfrak o}\) 的这样一组基。

\(\mathfrak B\) 包含在交 \([\mathfrak B_{\mathfrak o},\mathfrak O]\) 中。反过来，设 \(c\) 是这个收缩模的一个元素。作为 \(\mathfrak B_{\mathfrak o}\) 的元素，\(c\) 可以用 \(T\) 中的元素表示；同时，这也是用 \(\mathfrak O_{\mathfrak o}\) 的 \(\mathfrak o\)-基 \(T,\overline T\) 作出的一个表示：
\[
c=\gamma_1t_1+\cdots+\gamma_rt_r
\qquad\hbox{其中 }\gamma\hbox{ 属于 }\mathfrak o.
\]
作为 \(\mathfrak O\) 的元素，\(c\) 又有一个用 \(\mathfrak O\) 的 \(\mathfrak h\)-基 \(T,\overline T\) 作出的表示，即：
\[
c=h_{i_1}t_{i_1}+\cdots+h_{i_r}t_{i_r}+h_{j_1}\overline t_{j_1}+\cdots+h_{j_s}\overline t_{j_s}.
\]
由于 \(\mathfrak O\) 是 \(\mathfrak O_{\mathfrak o}\) 的子环，而 \(\mathfrak h\) 是 \(\mathfrak o\) 的子环，这同时也是用 \(\mathfrak O_{\mathfrak o}\) 的 \(\mathfrak o\)-基 \(T,\overline T\) 作出的表示。因此，这两个表示的对应系数必相等；各 \(\overline t\) 的系数消失，并且
\[
c=h_1t_1+\cdots+h_rt_r
\]
表明 \(c\) 是 \(\mathfrak B\) 的元素。由此，\(\mathfrak B\) 已被识别为收缩模。

\paragraph{附注。}
证明所用的假设比定理中表述的略弱，并且还证明了其余各项；因而有如下加强的表述：

\emph{如果 \(\mathfrak O\) 有一组独立 \(\mathfrak h\)-模基 \(T,\overline T\)，其中 \(T\) 包含于 \(\mathfrak B\)，并且 \(T\) 同时是 \(\mathfrak B_{\mathfrak o}\) 的一组（独立）\(\mathfrak o\)-模基，那么 \(\mathfrak B\) 是其扩张 \(\mathfrak B_{\mathfrak o}\) 的收缩模。此时还可证明 \(T\) 是 \(\mathfrak B\) 的一组独立 \(\mathfrak h\)-模基。}

\paragraph{3. 直积存在的一个充分判别准则。}
有时（例如在相对不同的情形中）只有在扩张系数域 \(\mathfrak h\) 之后，独立模基才存在。这时可由下列充分判别准则推得直积存在：

\emph{若存在 \(\mathfrak h\) 的一个扩张环 \(\mathfrak f\)，使得关于 \(\mathfrak h\) 的直积 \(\mathfrak O_{\mathfrak f}\) 和 \(\mathfrak o_{\mathfrak f}\) 存在，并且关于 \(\mathfrak f\) 的直积 \(\mathfrak O_{\mathfrak f}\times\mathfrak o_{\mathfrak f}\) 也存在，那么关于 \(\mathfrak h\) 的 \(\mathfrak O\times\mathfrak o\) 存在。}

须证明 \(\mathfrak R=\mathfrak O_{\mathfrak f}\times\mathfrak o_{\mathfrak f}\) 是 § 1, 2 意义下的嵌入环。事实上，\(\mathfrak O\) 与 \(\mathfrak o\) 作为 \(\mathfrak O_{\mathfrak f}\) 与 \(\mathfrak o_{\mathfrak f}\) 的子环，其元素两两可交换；又有
\[
[\mathfrak O,\mathfrak o]\subseteq[\mathfrak O_{\mathfrak f},\mathfrak o_{\mathfrak f}]=\mathfrak f;
\]
同样，\([\mathfrak O,\mathfrak o]\subseteq\mathfrak O\)；所以 \([\mathfrak O,\mathfrak o]\subseteq[\mathfrak O,\mathfrak f]=\mathfrak h\)，从而 \([\mathfrak O,\mathfrak o]=\mathfrak h\)。由此证明了 \(\mathfrak R\) 是一个嵌入环。

\paragraph{注。}
上述证明只对 \(\mathfrak O_{\mathfrak f}\) 使用了关于交的假设，对 \(\mathfrak o_{\mathfrak f}\) 并未使用；所以，后一个环也可以仅仅是 \(\mathfrak o\) 与 \(\mathfrak f\) 的乘积。
\subsection*{§ 3. 环相对于子环的不同。定义理想的微分商。}

从现在起，假定所出现的一切环都是交换环。\srcfn{12)}{若无这一假定，就必须定义右不同与左不同，并且必须考察二者彼此符合到何种程度，以及在半单系统的阶的情形下，它们同已知构造符合到何种程度（作为 § 4 中研究的推广）。}此外，始终假定环 $\mathfrak O$ 和 $\mathfrak o$ 是子环 $\mathfrak h$ 的同构且等价的扩张（亦即从 $\mathfrak O$ 到 $\mathfrak o$ 的映射是 $\mathfrak h$ 上恒等映射的延拓），并且关于 $\mathfrak h$ 的直积 $\mathfrak O_{\mathfrak o}$ 存在。

\paragraph{1. 不同的定义。}
假定关于 $\mathfrak h$ 的直积 $\mathfrak O_{\mathfrak o}$ 存在。令 $x$ 遍历 $\mathfrak O$ 的所有元素，并令 $\xi$ 每次表示 $\mathfrak o$ 中由该同构所对应的元素；以 $\mathfrak B$ 表示由所有差 $x-\xi$ 生成的 $\mathfrak O_{\mathfrak o}$ 的双边\emph{差分理想}，
\[
\mathfrak B=\{\ldots,(x-\xi),\ldots\}.\srcfnmark{13)}
\]
\srcfntext{13)}{这一记号将始终用于由集合生成的理想。}
\emph{差分商} $\mathfrak A$ 定义为零理想关于 $\mathfrak B$ 的理想商，因而由 $\mathfrak O_{\mathfrak o}$ 中所有使 $a\mathfrak B$ 消失的元素 $a$ 组成，
\[
\mathfrak A=(0):\mathfrak B.
\]
$\mathfrak o$ 关于 $\mathfrak h$ 的\emph{不同} $\mathfrak d$ 定义为 $\mathfrak A[x\to\xi]$，也就是说，它由 $\mathfrak A$ 产生，方法是在 $\mathfrak A$ 的每个 $a=\alpha_1x_1+\cdots+\alpha_rx_r$ 中，将各 $x$ 换成与之同构对应的 $\xi$。相应地，$\mathfrak O$ 关于 $\mathfrak h$ 的不同 $\mathfrak D$ 定义为 $\mathfrak A[\xi\to x]$；由于 $\mathfrak A$ 的定义关于 $x$ 与 $\xi$ 对称，$\mathfrak d$ 与 $\mathfrak D$ 是由交换 $\xi$ 与 $x$ 得到的，因此在从 $\mathfrak o$ 到 $\mathfrak O$ 的同构下彼此对应。由于 $\mathfrak A$ 是 $\mathfrak O_{\mathfrak o}$ 中的理想，特别地既是 $\mathfrak o$-模又是 $\mathfrak O$-模，故 $\mathfrak d$ 与 $\mathfrak D$ 分别成为 $\mathfrak o$ 与 $\mathfrak O$ 中的理想。

\paragraph{2. 定义理想的微分商。}
设（§ 1, 3）$\mathfrak M$ 是 $\mathfrak O$ 关于 $\mathfrak h$ 的定义理想，与 $z$ 的生成系相对应。于是（§ 1, 3）$\mathfrak M$ 也是 $\mathfrak o$ 关于 $\mathfrak h$ 的定义理想，与 $\xi$ 的同构生成系相对应；而且（§ 1, 4）$\mathfrak O_{\mathfrak o}$ 同构于 $\mathfrak o[\ldots Z\ldots]/\mathfrak M_{\mathfrak o}$。现在，差分理想与差分商在多项式环 $\mathfrak o[\ldots Z\ldots]$ 中分别定义为
\[
\mathfrak Q=\{\ldots,(Z-\xi),\ldots\}
\]
以及
\[
\mathfrak C=\mathfrak M_{\mathfrak o}:\mathfrak Q
=\{\ldots(F(Z)-F(\xi)),\ldots\}:\{\ldots,(Z-\xi),\ldots\}.
\]
\emph{微分商} $\mathfrak M'$ 定义为 $\mathfrak C[\xi\to Z]$。这里应当注意，这一定义是唯一的，并不依赖于用生成元 $\xi$ \emph{表示} $\mathfrak o$ 中元素的方式。事实上，由于 $\mathfrak C$ 是理想商，它包含 $\mathfrak M_{\mathfrak o}$；同一个 $\mathfrak o$ 中元素的两个表示相差一个多项式 $F(\xi)$，因而在代入 $Z$ 后，相差一个属于 $\mathfrak M$、从而已经属于 $\mathfrak C$ 的 $F(Z)$。理想 $\mathfrak M'$ 成为 $\mathfrak h[\ldots Z\ldots]$ 中包含 $\mathfrak M$ 的理想。\emph{微分商 $\mathfrak M'$ 在 $Z\to\xi$ 时过渡为 $\mathfrak o$ 关于 $\mathfrak h$ 的不同 $\mathfrak d$。}这是因为在同态映射下，不仅理想彼此对应，商也彼此对应。因此，在 $\mathfrak o[\ldots Z\ldots]$ 与 $\mathfrak O_{\mathfrak o}$ 之间、即使 $Z$ 与 $z$ 对应的同态下，$\mathfrak A$ 成为 $\mathfrak C$ 的同态像，因为相对于 $\mathfrak Q$ 与 $\mathfrak M_{\mathfrak o}$，$\mathfrak B$ 与零理想也有同样的关系。\srcfn{13a)}{编者注：为了能够断定商彼此对应，关键在于 $\mathfrak M_{\mathfrak o}$ 表示 $(0)$ 的原像集合。}由此可知，在该同态下，$\mathfrak C$ 中的代换 $\xi\to Z$ 对应于 $\mathfrak A$ 中的代换 $\xi\to z$；因此，在 $Z$ 与 $z$ 对应之下，$\mathfrak O$ 的不同 $\mathfrak D$ 成为 $\mathfrak M'$ 的同态像。所以
\[
\mathfrak M'[Z\to z]=\mathfrak D
\qquad\hbox{且}\qquad
\mathfrak M'[Z\to\xi]=\mathfrak D[z\to\xi]=\mathfrak d.
\]

\paragraph{3. 有一个定义方程时的不同。}
设 $\mathfrak O$ 有一个定义方程 $G(z)=0$（§ 1, 3 末尾），并进一步\emph{假定}最高次项的系数等于单位元，即
\[
G(Z)=Z^n+h_1Z^{n-1}+\cdots+h_n
\]
且系数 $h_i$ 属于 $\mathfrak h$。那么，它与多项式微分商之间的联系由下述结果给出：

\paragraph{定理。}
\emph{$\mathfrak o$ 或 $\mathfrak O$ 关于 $\mathfrak h$ 的不同有定义（即关于 $\mathfrak h$ 的直积 $\mathfrak O_{\mathfrak o}$ 存在），并且分别成为以 $G'(\xi)$ 或 $G'(z)$ 为生成元的主理想，其中 $G'(Z)$ 表示多项式微分商。理想微分商 $\mathfrak M'$ 等于 $(G'(Z),G(Z))$。}

\paragraph{证明。}
由假定首先可知，元素 $e,z,\ldots,z^{n-1}$ 构成 $\mathfrak O$ 的一个独立模基。事实上，因为 $G(Z)$ 的最高次系数是单位元，这些幂确实构成一个 $\mathfrak h$-模基；而这个模基是独立的，因为 $\mathfrak M$ 不可能包含次数低于 $n$ 的多项式。多项式 $K(Z)G(Z)$ 的次数确实等于两个因子次数之和，这仍然是由关于最高次系数的假定所致。于是，由于该模基存在，关于 $\mathfrak h$ 的直积 $\mathfrak O_{\mathfrak o}$ 存在，从而不同有定义。其余结论由考察 $\mathfrak o[\ldots Z\ldots]$ 中的差分商 $\mathfrak C$ 得到（参见 2）。这个差分商成为以如下多项式为生成元的主理想：
\[
\begin{aligned}
C(Z,\xi)={}&(Z^{n-1}+Z^{n-2}\xi+\cdots+\xi^{n-1})\\
&+h_1(Z^{n-2}+\cdots+\xi^{n-2})+\cdots+h_{n-1}.
\end{aligned}
\]
事实上，$C$ 由 $G(Z)-G(\xi)$ 除以 $Z-\xi$ 的形式除法得到，因而属于 $\mathfrak C$，因为 $Z-\xi$ 成为 $\mathfrak Q$ 的生成元。任意多项式模 $C$ 均可约化为一个关于 $Z$ 的次数达不到 $n-1$ 的多项式 $D$；于是乘积 $D(Z-\xi)$ 的次数达不到 $n$，所以只有当它消失时才可能属于 $\mathfrak M$。但是，由 $D(Z-\xi)=0$ 可推出 $D=0$，因为 $Z-\xi$ 中 $Z$ 的系数是单位元。所以 $C$ 成为 $\mathfrak C$ 的生成元。而由在 $C$ 及其倍式中代换 $\xi\to Z$（亦即 $G(Z)=C(Z-\xi)$ 及其倍式），现在得到
\[
\mathfrak M'=(G'(Z),G(Z)),
\]
进而得到
\[
\mathfrak d=(G'(\xi))
\qquad\hbox{且}\qquad
\mathfrak D=(G'(z)).
\]
% Source: Paper 43, printed pp. 10--13 (journal scan pp. 699--702)
\subsection*{\S{} 4. 直和的不同。}

对 $\mathfrak O$ 作如下假定：

\begin{enumerate}
\item[(1)] $\mathfrak O$ 是诸理想的直和
\[
\mathfrak O=\mathfrak R_1+\cdots+\mathfrak R_r
=\mathfrak R_i+\mathfrak S_i
=\mathfrak Oe_1+\cdots+\mathfrak Oe_r,
\]
其中 $e_i$ 表示单位元的各个分量。

\item[(2)] 每个 $\mathfrak R_i$ 都有一个独立的 $\mathfrak h_i$-模基 $T_i$，其中 $\mathfrak h_i$ 是 $\mathfrak h$ 在 $\mathfrak R_i$ 中的分量，即 $\mathfrak h_i=\mathfrak h e_i$。这个基 $T_i$ 包含 $\mathfrak R_i$ 的单位元 $e_i$ 作为一个元素。

\item[(3)] 每个交集 $[\mathfrak h,\mathfrak S_i]$ 都等于零元（$\mathfrak S_i$ 不含常数）。
\end{enumerate}

由于存在同构，同样的假定对 $\mathfrak o$ 也成立。

\paragraph{定理。}
\emph{在上述假定下，$\mathfrak O$ 关于 $\mathfrak h$ 的不同成为各 $\mathfrak R_i$ 关于 $\mathfrak h_i$ 的不同的直和。}

证明以证实一系列个别事实为基础。

\paragraph{I.}
首先可知，$\mathfrak O$ 也有一个包含 $e$ 的独立 $\mathfrak h$-模基，因而 $\mathfrak O$ 或 $\mathfrak o$ 关于 $\mathfrak h$ 的不同有定义；具体而言，这个基由 $\mathfrak R_i$ 的各个基 $T_i$ 的并集 $T$ 构成。事实上，由
\[
x=xe_1+\cdots+xe_r
\]
以及
\[
xe_i=(h_1e_i)t_{1i}e_i+\cdots+(h_\lambda e_i)t_{\lambda i}e_i
=h_1(t_{1i}e_i)+\cdots+h_\lambda(t_{\lambda i}e_i)
\]
可知，$T$ 确实成为 $\mathfrak O$ 的一个 $\mathfrak h$-基。这个基还是独立的；因为由 $T$ 中元素之间的一个关系，借助直和可得到每个 $T_i$ 中元素之间的关系，因而对所有系数都有 $h e_i=0$。但根据假定 (3)，这又给出 $h=0$；事实上，由于
\[
h\equiv h e_i\pmod{\mathfrak S_i}
\]
故从 $h e_i=0$ 可知 $h$ 属于 $[\mathfrak h,\mathfrak S_i]$。此外，$T$ 包含 $\mathfrak R_i$ 的各个单位元 $e_i$，而它们的和是 $e$，所以可在 $T$ 中以 $e$ 替换一个 $e_i$，由此得到满足一切条件的 $\mathfrak h$-基 $\overline T$。于是便得到 $\mathfrak O_{\mathfrak o}$ 以及 $\mathfrak O$ 关于 $\mathfrak h$ 的不同的存在性。

\paragraph{II.}
交集条件 (3) 现在蕴含如下更强的条件：
\[
[\mathfrak o,\mathfrak S_{i(\mathfrak o)}]=0,
\]
其中 $\mathfrak S_{i(\mathfrak o)}$ 表示 $\mathfrak O_{\mathfrak o}$ 中由 $\mathfrak S_i$ 导出的理想。设
\[
\alpha=\alpha e=\beta_1t_1+\cdots+\beta_\lambda t_\lambda,
\]
其中 $t_1,\ldots,t_\lambda$ 是从 $T$ 中删去 $T_i$ 所得的基的元素，而 $\alpha$ 与 $\beta$ 属于 $\mathfrak o$。若以 $e$ 替换 $e_i$ 而构成 $\overline T$，那么 $e,t_1,\ldots,t_\lambda$ 就是 $\mathfrak O$ 的 $\mathfrak h$-基 $\overline T$ 中彼此不同的元素。但是，由于 $\mathfrak O_{\mathfrak o}$ 中的等式在每个这种基中都化为系数相等，关系
\[
\alpha=\alpha e=\beta_1t_1+\cdots+\beta_\lambda t_\lambda
\]
便给出 $\alpha=0$，从而\emph{证明了加强的交集性质}。由此特别得到（如上）：由
\[
\alpha\text{ 是 }\mathfrak o\text{ 的元素}\quad\text{且}\quad\alpha\ne0
\]
可推出 $\alpha e_i\ne0$。

\paragraph{III.}
这个交集条件的一个直接推论是：$\mathfrak o$ 与它的、对应于 $\mathfrak O$ 的分解的各个\emph{分量} $\mathfrak o e_i$ \emph{环同构}；此同构包含同构
\[
\mathfrak h\simeq\mathfrak h_i=\mathfrak h e_i.
\]
事实上，由于 $e_i^2=e_i$，并且 $e_i$ 与 $\mathfrak o$ 的所有元素可交换，$\mathfrak o e_i$ 成为 $\mathfrak o$ 的环同态像；这一对应又是一一的，因为由 $\alpha\ne0$ 也可推出 $\alpha e_i\ne0$。在这一对应下，子环 $\mathfrak h$ 对应于分量 $\mathfrak h_i=\mathfrak h e_i$。

\paragraph{IV.}
若进一步采用 $\mathfrak o$ 的相应分解，便得到\emph{另外的同构}：
\[
\mathfrak o=\mathfrak r_1+\cdots+\mathfrak r_r
=\mathfrak r_i+\mathfrak s_i
=\mathfrak o\varepsilon_1+\cdots+\mathfrak o\varepsilon_r.
\]
于是有：
\[
\mathfrak o\simeq\mathfrak o e_i,
\quad \mathfrak O\simeq\mathfrak O\varepsilon_i,
\quad \mathfrak r_i\simeq\mathfrak r_i e_i,
\quad \mathfrak R_i\simeq\mathfrak R_i\varepsilon_i;
\quad \mathfrak h\simeq\mathfrak h e_i\simeq\mathfrak h\varepsilon_i\simeq\mathfrak h e_i\varepsilon_i.
\]
如果在上述考察中从 $\mathfrak o$ 而不是 $\mathfrak O$ 出发，关系 $\mathfrak O\simeq\mathfrak O\varepsilon_i$ 就与 $\mathfrak o\simeq\mathfrak o e_i$ 完全对应。同样，接下来的两个关系彼此对应；其中第一个关系来自如下事实：在 $\mathfrak o\simeq\mathfrak o e_i$ 下，子环 $\mathfrak r_i$ 也对应于 $\mathfrak r_i e_i$；同理得到
\[
\mathfrak h\simeq\mathfrak h e_i\simeq\mathfrak h\varepsilon_i.
\]
而由于 $\mathfrak h\varepsilon_i$ 是 $\mathfrak o$（也是 $\mathfrak r_i$）的子环，最后还有
\[
\mathfrak h\varepsilon_i\simeq\mathfrak h e_i\varepsilon_i.
\]

\paragraph{V.}
由 $\mathfrak O$ 的直和表示可知：\emph{若 $\mathfrak C=(0):\mathfrak B$ 是 $\mathfrak O$ 中零理想关于理想 $\mathfrak B$ 的商，则每个分量 $\mathfrak C_i$ 等于 $\mathfrak R_i$ 的零理想关于分量 $\mathfrak B_i$ 的商。}由
\[
\mathfrak B\mathfrak C=(\mathfrak B e_1)(\mathfrak C e_1)+\cdots+(\mathfrak B e_r)(\mathfrak C e_r)=0
\]
根据\emph{直和}可逐一推出
\[
(\mathfrak B e_i)(\mathfrak C e_i)=0.
\]
反过来，若 $\mathfrak R_i$ 中的元素 $r e_i$ 满足条件 $(\mathfrak B e_i)(r e_i)=0$，那么对 $i\ne k$ 也有 $(\mathfrak B e_k)(r e_i)=0$。

所以 $r e_i$ 是 $\mathfrak C$ 的元素，同时也是 $\mathfrak R_i$ 的元素，因而属于分量 $\mathfrak C_i=\mathfrak C e_i$。由此识别出 $\mathfrak C_i=\mathfrak C e_i$ 是零理想关于 $\mathfrak B_i$ 的商；它也就是 $\mathfrak R_i$ 的零理想关于 $\mathfrak B_i$ 的商，因为由
\[
r e_i\mathfrak B_i\equiv0\quad(\mathfrak S_i)
\]
总能推出 $r e_i\mathfrak B_i=0$。\srcfn{14)}{本定理也可直接由第 2 节所用的定理得到，即在同态映射下，不仅理想彼此对应，商也彼此对应。}由于除 $\mathfrak O$ 外，
\[
\mathfrak O_{\mathfrak o}=\mathfrak O_{\mathfrak o}e_1+\cdots+\mathfrak O_{\mathfrak o}e_r
\]
也成为直和，\emph{所以该定理对 $\mathfrak O_{\mathfrak o}$ 中的理想也成立。}

\paragraph{VI. 不同的分量分解。}
按定义（参见 § 3, 1），原有 $\mathfrak D=\mathfrak A[\xi\to x]$；故有
\[
\mathfrak D=\mathfrak D e_1+\cdots+\mathfrak D e_r
=\mathfrak A[\xi\to x]e_1+\cdots+\mathfrak A[\xi\to x]e_r.
\]
但是，由于 $e_i$ 属于 $\mathfrak O$，因而对应于其自身；又由于 $\varepsilon_i$ 在该同构下过渡为 $e_i$，并且 $e_i\varepsilon_i=e_i$，所以
\[
\mathfrak A[\xi\to x]e_i=\mathfrak A e_i[\xi\to x]
=\mathfrak A e_i\varepsilon_i[\xi\to x].
\]
而根据 V，$\mathfrak A e_i\varepsilon_i$ 等于商 $(0):\mathfrak B e_i\varepsilon_i$，也等于 $\mathfrak R_i\mathfrak r_i$ 的零理想关于 $\mathfrak B e_i\varepsilon_i$ 的商。事实上，$\mathfrak O$ 与 $\mathfrak o$ 的直和分解经乘法产生分解
\[
\mathfrak O_{\mathfrak o}=\sum\mathfrak O_{\mathfrak o}e_i\varepsilon_k
=\sum\mathfrak R_i\mathfrak r_k,
\]
这个分解是直和，因为 $\sum e_i\varepsilon_k=e$，并且这些单位元满足正交关系。

\paragraph{VII.}
\emph{证明 $\mathfrak D e_i=\mathfrak A e_i\varepsilon_i[\xi\to x]=(0):\mathfrak B e_i\varepsilon_i$ 是 $\mathfrak R_i$ 关于 $\mathfrak h_i$ 的不同。}这一证明基于下列事实：$\mathfrak R_i\mathfrak r_i$ 等于关于 $\mathfrak h e_i\varepsilon_i$ 的直积 $\mathfrak R_i\varepsilon_i\times\mathfrak r_i e_i$；其差分理想由 $\mathfrak B e_i\varepsilon_i$ 给出，其差分商由 $\mathfrak A e_i\varepsilon_i$ 给出，而 $\mathfrak D e_i\varepsilon_i$ 成为 $\mathfrak R_i\varepsilon_i$ 的不同，相应地 $\mathfrak d e_i\varepsilon_i$ 成为 $\mathfrak r_i e_i$ 的不同。于是，同构
\[
\mathfrak R_i\varepsilon_i\simeq\mathfrak R_i
\qquad\text{和}\qquad
\mathfrak r_i e_i\simeq\mathfrak r_i
\]
分别给出 $\mathfrak D e_i$ 与 $\mathfrak d\varepsilon_i$，作为 $\mathfrak R_i$ 关于 $\mathfrak h e_i$ 的不同与 $\mathfrak r_i$ 关于 $\mathfrak h\varepsilon_i$ 的不同。

\paragraph{VIII.}
\emph{$\mathfrak R_i\mathfrak r_i$ 等于关于 $\mathfrak h e_i\varepsilon_i$ 的直积 $\mathfrak R_i\varepsilon_i\times\mathfrak r_i e_i$。}由于 $\mathfrak R_i$ 有 $\mathfrak h_i$-模基 $T_i$，从同构 (IV) $\mathfrak R_i\simeq\mathfrak R_i\varepsilon_i$ 可知，$T_i\varepsilon_i$ 成为 $\mathfrak R_i\varepsilon_i$ 的一个 $\mathfrak h e_i\varepsilon_i$-模基。现在，根据 IV，$\mathfrak r_i e_i$ 同构于 $\mathfrak R_i\varepsilon_i$——因为借助 $\mathfrak o\simeq\mathfrak O$ 有 $\mathfrak r_i\simeq\mathfrak R_i$——并且是 $\mathfrak h e_i\varepsilon_i$ 的扩张环。因此，该直积由 $T_i\varepsilon_i$ 中元素的一切线性形式组成，系数取自 $\mathfrak r_i e_i$，而等式定义为各系数相等。但借助 $\mathfrak O\times\mathfrak o$，这恰好成为 $\mathfrak R_i\mathfrak r_i$，因为 $T_i=T_i e_i$、$\mathfrak r_i=\mathfrak r_i\varepsilon_i$，并且 $T_i$ 属于 $\mathfrak O$ 的一个模基。

\paragraph{IX.}
\emph{$\mathfrak R_i\varepsilon_i\times\mathfrak r_i e_i$ 的差分理想由 $\mathfrak B e_i\varepsilon_i$ 给出，差分商由 $\mathfrak A e_i\varepsilon_i$ 给出。}为看清这一点，把 $\mathfrak B$ 写成如下形式：
\[
\mathfrak B=\{\ldots,xe_1-\xi\varepsilon_1,\ldots,xe_r-\xi\varepsilon_r,\ldots\}.
\]
于是
\[
\mathfrak B e_i\varepsilon_i
=\{\ldots,xe_i\varepsilon_i-\xi e_i\varepsilon_i,\ldots\},
\]
所以它确实等于 $\mathfrak R_i\varepsilon_i\times\mathfrak r_i e_i$ 的差分理想。而因为（根据 VI）$\mathfrak A e_i\varepsilon_i$ 等于此环的零理想关于 $\mathfrak B e_i\varepsilon_i$ 的商，所以它被识别为差分商。

\paragraph{X.}
\emph{$\mathfrak R_i\varepsilon_i$ 的不同等于 $\mathfrak D e_i\varepsilon_i$。}$\mathfrak R_i\varepsilon_i$ 的不同通过如下方式产生：在 $\mathfrak A e_i\varepsilon_i$ 中，以每个 $\mathfrak r_i e_i$ 中元素的同构对应元素替换它；也就是每次以 $(xe_i)\varepsilon_i$ 替换 $(\xi\varepsilon_i)e_i$。换言之，由于 $\xi\varepsilon_i$ 也是 $\mathfrak o$ 中的一个 $\xi$，它过渡为 $xe_i$：在 $\mathfrak A e_i\varepsilon_i$ 中作代换 $\xi\to x$——这把 $(\xi\varepsilon_i)e_i$ 变为 $(xe_i)e_i=xe_i$——然后乘以 $\varepsilon_i$。而（根据 VI）原有 $\mathfrak A e_i\varepsilon_i[\xi\to x]=\mathfrak D e_i$，所以 $\mathfrak R_i\varepsilon_i$ 关于 $\mathfrak h e_i\varepsilon_i$ 的不同等于 $\mathfrak D e_i\varepsilon_i$；与此对称，$\mathfrak r_i e_i$ 的不同等于 $\mathfrak d e_i\varepsilon_i$。

\paragraph{XI.}
\emph{$\mathfrak R_i$ 的不同等于 $\mathfrak D e_i$。}事实上，在该同构下，$\mathfrak R_i$ 的不同通过乘以 $\varepsilon_i$ 过渡为 $\mathfrak R_i\varepsilon_i$ 的不同；而后者等于 $\mathfrak D e_i\varepsilon_i$，所以前者等于 $\mathfrak D e_i$。因此：
\[
\mathfrak D=\mathfrak D e_1+\cdots+\mathfrak D e_r,
\]
\emph{并且 $\mathfrak D e_i$ 是 $\mathfrak R_i$ 的不同；}与此对称，
\[
\mathfrak d=\mathfrak d\varepsilon_1+\cdots+\mathfrak d\varepsilon_r,
\]
\emph{并且 $\mathfrak d\varepsilon_i$ 是 $\mathfrak r_i$ 的不同，}定理得证。
% Source: Paper 43, printed pp. 13--17 (journal scan pp. 702--706)
\subsection*{\S{} 5. 域的伽罗瓦扩张环。结构定理。}

\paragraph{1. 基础环。}
设 $K$ 是一个\emph{基域 $P$ 上的 $n$ 次域}，并且假定 $K$ 是可分扩张。设 $\Gamma$ 是相应的伽罗瓦域；它在同构意义下唯一确定，包含 $n$ 个共轭域 $K^{(i)}$，并且是它们的直积，其中规定 $K=K^{(1)}$。用 $\mathfrak K$ 表示 $P$ 的一个扩张，它同 $K$、因而也同所有 $K^{(i)}$ 同构，并且相对于 $P$ 与它们等价；此外还要求其与 $\Gamma$ 的交 $[\mathfrak K,\Gamma]$ 就是 $P$（亦即 $\Gamma$ 与 $\mathfrak K$ 不存在共同的包络域）。\srcfn{15)}{这里当然假定——这并不限制一般性——$K$ 与 $\mathfrak K$ 中不属于 $P$ 的元素之间不存在任何关系（参见第 1 节第 2 小节、第 2 节第 1 小节以及脚注 11）。}由于 $\mathfrak K$ 具有一个独立、有限且包含单位元的 $P$-模基，根据第 2 节第 1 小节，直积 $\mathfrak K_\Gamma$、$\mathfrak K_K$、$\mathfrak K_{K^{(i)}}$ 都存在；$\mathfrak K_\Gamma$ 是它的 $n$ 个子环 $\mathfrak K_{K^{(i)}}$ 的直积。称 $\mathfrak K_\Gamma$ 为 $\mathfrak K$ 的\emph{伽罗瓦扩张环}。同时，$\mathfrak K_\Gamma$ 还是 $\mathfrak K_K$ 与 $\Gamma$ 关于 $K$ 的直积，因为 $\mathfrak K$ 的每个独立 $P$-基也成为 $\mathfrak K_K$ 的独立 $K$-基；对各共轭域也同样如此。\srcfn{16)}{对所有共轭域一致成立的定理，通常只对 $\mathfrak K_K$ 表述，也就是省去指标。}这些环的理想之间有如下关系：

\paragraph{定理。}
\emph{$\mathfrak K$ 中的每个理想，都是它在 $\mathfrak K_K$ 或 $\mathfrak K_\Gamma$ 中的扩张的收缩理想；$\mathfrak K_K$ 中的每个理想，都是它在 $\mathfrak K_\Gamma$ 中的扩张的收缩理想。}

\paragraph{证明。}
由于 $P$ 和 $K$ 都是域，$\mathfrak K$ 中的每个 $P$-模以及 $\mathfrak K_K$ 中的每个 $K$-模，都分别具有一个独立有限的 $P$-模基或 $K$-模基，而且该基可以分别补充成 $\mathfrak K$ 或 $\mathfrak K_K$ 的这种基。于是该定理由第 2 节第 2 小节中的定理推出。

用 $\mathfrak B_{K^{(i)}}^{(i)}$ 表示 $\mathfrak K_{K^{(i)}}$ 的差分理想，即
\[
\mathfrak B_{K^{(i)}}^{(i)}=\{\ldots,x-\xi^{(i)},\ldots\},
\]
特别地：
\[
\mathfrak B_K=\{\ldots,x-\xi,\ldots\}.
\]
由于 $\xi$ 或 $\xi^{(i)}$ 分别是 $K$ 或 $K^{(i)}$ 中通过同构与 $x$ 对应的元素，所以这 $n$ 个差分理想彼此共轭。

用 $\mathfrak A_K$ 或 $\mathfrak A_{K^{(i)}}^{(i)}$ 表示 $\mathfrak K_K$ 中的差分商 $(0):\mathfrak B_K$ 及其各共轭。它们在 $\mathfrak K_\Gamma$ 中的扩张理想分别记作 $\mathfrak B_\Gamma$ 和 $\mathfrak A_\Gamma$。由本小节关于扩张理想与收缩理想的定理，特别得到：

\emph{差分理想 $\mathfrak B_K$ 和差分商 $\mathfrak A_K$，分别是它们的扩张 $\mathfrak B_\Gamma$ 和 $\mathfrak A_\Gamma$ 的收缩理想。}

\paragraph{2. 伽罗瓦扩张环的直和分解。}
差分理想及差分商同伽罗瓦扩张环结构之间的关系，由下述定理给出。

\paragraph{结构定理。}
\emph{伽罗瓦扩张环 $\mathfrak K_\Gamma$ 的零理想，是 $n$ 个共轭差分理想的扩张之交；$\mathfrak K_\Gamma$ 本身，是 $n$ 个共轭差分商的扩张之直和。$\mathfrak K_{K^{(i)}}$ 是相应差分理想与差分商的直和，而它的零理想是二者之交。}

\paragraph{证明。}
按每个 $\mathfrak B_\Gamma^{(i)}$ 取的剩余类环关于 $\Gamma$ 的秩为一，因为每个元素模 $\mathfrak B_\Gamma^{(i)}$ 都同 $\Gamma$ 中的一个元素同余，而 $\mathfrak B_\Gamma^{(i)}$ 不可能包含 $\Gamma$ 中任何非零元素（因为令 $x\to\xi^{(i)}$ 时 $\mathfrak B_\Gamma^{(i)}$ 消失）。由于可分性，各 $\mathfrak B_\Gamma^{(i)}$ 全都不同，因而它们两两互素（任意两个之和等于 $\mathfrak K_\Gamma$）。因此，理想
\[
\mathfrak B_\Gamma^{(1)},
[\mathfrak B_\Gamma^{(1)},\mathfrak B_\Gamma^{(2)}],\ldots,
[\mathfrak B_\Gamma^{(1)},\ldots,\mathfrak B_\Gamma^{(n)}]
\]
的秩每一步至少降低一个单位；所有这些理想之交成为零理想，而由于两两互素，少于全体共轭理想的交必定异于零理想（所以秩在每一步恰好降低一个单位）。由熟知的定理，这便给出 $\mathfrak K_\Gamma$ 如下唯一确定的直和表示：\srcfn{17)}{例如参见 E. Noether，\emph{Abstrakter Aufbau der Idealtheorie in algebraischen Zahl- und Funktionenkörpern}，Math. Ann. 96 (1927)，26--61，第 4 节第 5 小节。}
\begin{align}
\mathfrak K_\Gamma
&=\mathfrak C_\Gamma^{(1)}+\cdots+\mathfrak C_\Gamma^{(n)}
=\Gamma e^{(1)}+\cdots+\Gamma e^{(n)}
=\mathfrak C_\Gamma^{(i)}+\mathfrak B_\Gamma^{(i)}, \tag{1}\\
e&=e^{(1)}+\cdots+e^{(n)},\qquad
\mathfrak C_\Gamma^{(i)}=
[\mathfrak B_\Gamma^{(1)},\ldots,\mathfrak B_\Gamma^{(i-1)},
\mathfrak B_\Gamma^{(i+1)},\ldots,\mathfrak B_\Gamma^{(n)}].\notag
\end{align}
还须证明 $\mathfrak C_\Gamma$ 等于 $\mathfrak A_\Gamma$。只要证明 $\mathfrak K_K$ 中相应的二者相等，这一结论便随之成立。在这里有分解：
\begin{equation}
\mathfrak K_K=\mathfrak C_K+\mathfrak B_K
=\mathfrak K_K e^{(1)}+\mathfrak K_K(e-e^{(1)}). \tag{2}
\end{equation}
事实上，$e^{(1)}$ 属于 $\mathfrak K_K$；因为 $\mathfrak C_\Gamma^{(1)}$ 在 $\mathfrak B_\Gamma^{(2)},\ldots,\mathfrak B_\Gamma^{(n)}$ 的所有置换下保持不变，所以作为 $\mathfrak C_\Gamma^{(1)}$ 中单位元分量的 $e^{(1)}$ 也同样保持不变。因此，用 $\mathfrak K_K$ 的一个独立 $K$-基来表示 $e^{(1)}$——这是可能的，因为该基同时也是 $\mathfrak K_\Gamma$ 的独立 $\Gamma$-基——则各个系数在 $K^{(2)},\ldots,K^{(n)}$ 的所有置换下都保持不变，因而属于 $K$。由于 $(e^{(1)})^2=e^{(1)}$，和式表示
\[
\mathfrak K_K=\mathfrak K_K e^{(1)}+\mathfrak K_K(e-e^{(1)})
\]
是直和；它在 $\mathfrak K_\Gamma$ 中的扩张仍然是直和，故由 (1) 得到 $\mathfrak K_\Gamma=\mathfrak C_\Gamma^{(1)}+\mathfrak B_\Gamma^{(1)}$。于是根据第 2 节第 2 小节的末尾部分，有
\[
\mathfrak K_K e^{(1)}=[\mathfrak C_\Gamma,\mathfrak K_K]
=[\Gamma e^{(1)},\mathfrak K_K]=Ke^{(1)}=\mathfrak C_K;
\]
\[
\mathfrak K_K(e-e^{(1)})=[\mathfrak B_\Gamma,\mathfrak K_K]=\mathfrak B_K,
\]
从而 (2) 得证。

由 (2) 又可推出 $\mathfrak C_K=\mathfrak A_K$，进而 $\mathfrak C_\Gamma=\mathfrak A_\Gamma$。事实上，由 $e^{(1)}(e-e^{(1)})=0$ 可得 $\mathfrak C_K\mathfrak B_K=0$；反过来，$b\mathfrak B_K=0$ 也给出 $b(e-e^{(1)})=0$，所以 $b=be=be^{(1)}$，于是
\[
\mathfrak C_K=(0):\mathfrak B_K=\mathfrak A_K,
\]
并且扩张后有 $\mathfrak C_\Gamma=\mathfrak A_\Gamma$。因此：
\begin{align}
\mathfrak K_\Gamma
&=\mathfrak A_\Gamma^{(1)}+\cdots+\mathfrak A_\Gamma^{(n)}
=\Gamma e^{(1)}+\cdots+\Gamma e^{(n)};
&0&=[\mathfrak B_\Gamma^{(1)},\ldots,\mathfrak B_\Gamma^{(n)}]; \tag{3}\\
\mathfrak K_{K^{(i)}}
&=\mathfrak A_{K^{(i)}}^{(i)}+\mathfrak B_{K^{(i)}}^{(i)}
=K^{(i)}e^{(i)}+\mathfrak B_{K^{(i)}}^{(i)};
&0&=[\mathfrak A_{K^{(i)}}^{(i)},\mathfrak B_{K^{(i)}}^{(i)}].\notag
\end{align}
至此，该定理的各部分全都得证。还应指出，这是一个\emph{完全不变的结构定理}，因为其中出现的所有构造（$\mathfrak K_K$、$\mathfrak K_\Gamma$、$\mathfrak B_K$、$\mathfrak A_K$）都不依赖于基来定义。\srcfn{18)}{这里也可以借助定义方程以及相应多项式的线性因式分解来表述，但这样的表述并非不变的。}

\paragraph{3. 由共轭元素给出的 $\mathfrak K$ 的分量表示。}
与共轭差分理想及差分商相对应，$\mathfrak K$ 的元素可通过共轭分量来表示：

\paragraph{定理。}
\emph{设
\[
x=\xi^{(1)}e^{(1)}+\cdots+\xi^{(n)}e^{(n)}
\]
是 $\mathfrak K$ 中元素 $x$ 的 $\mathfrak K$-分量表示，则这 $n$ 个分量彼此共轭。特别地，$\xi^{(1)},\ldots,\xi^{(n)}$ 正是通过 $\mathfrak K$ 到 $K^{(1)},\ldots,K^{(n)}$ 的各同构与 $x$ 对应的 $n$ 个共轭值。}

$e^{(1)},\ldots,e^{(n)}$ 彼此共轭，可由第 2 小节中的关系式 (3) 看出：把 $\mathfrak B_K$ 映到 $\mathfrak B_{K^{(i)}}^{(i)}$ 的 $\Gamma$ 的自同构，也同时把 $e^{(1)}$ 映到 $e^{(i)}$，而后者确实是 $\mathfrak K_{K^{(i)}}$ 中的元素。同一组关系式还给出 $xe^{(i)}=\xi^{(i)}e^{(i)}$；这是因为 $\mathfrak B_{K^{(i)}}^{(i)}e^{(i)}=0$，从而
\[
(x-\xi^{(i)})e^{(i)}=0
\]
（也可由以下事实推出：作为域的 $\mathfrak K$ 同它的分量同构，而该分量就是 $K^{(i)}e^{(i)}$，其中 $x$ 与 $\xi^{(i)}e^{(i)}$ 同构对应）。

由该定理可得：

\paragraph{推论。}
\emph{设 $\mathfrak B$ 是 $\mathfrak K$ 中的一个绝对模（即它在包含任意两个元素的同时也包含二者之差），并设 $\mathfrak B^{(1)},\ldots,\mathfrak B^{(n)}$ 是它在 $\mathfrak K_\Gamma$ 中的各分量，则 $\mathfrak B$ 等于交 $[\mathfrak K,\mathfrak B^{(1)}+\cdots+\mathfrak B^{(n)}]$。}

事实上，$\mathfrak B$ 包含在这个交中；反过来，若 $\mathfrak B^{(1)}+\cdots+\mathfrak B^{(n)}$ 中的一个元素还属于 $\mathfrak K$，则根据该定理，它必为共轭分量之和，因而是 $\mathfrak B$ 中的元素。

\paragraph{4. $\mathfrak K$ 的互补基及其同单位元分量的关系。}
从不变结构转到基表示时，这一结构把各个基归并为互补基对：

\paragraph{定理。}
\emph{$\mathfrak K$ 的每个 $P$-基 $t_1,\ldots,t_n$，都对应着 $\mathfrak K$ 中的一个互补基 $T_1,\ldots,T_n$，而后者的互补基又是原来的 $t$ 基。$K^{(i)}$ 中同互补基 $T_1,\ldots,T_n$ 同构的基 $A_1^{(i)},\ldots,A_n^{(i)}$，被定义为单位元分量 $e^{(i)}$ 在由 $t$ 给出的基表示
\[
e^{(i)}=A_1^{(i)}t_1+\cdots+A_n^{(i)}t_n
\]
中的系数。若 $\alpha_1^{(i)},\ldots,\alpha_n^{(i)}$ 是与各 $t_i$ 对应的基，则 $\alpha$ 与 $A$ 的矩阵互逆。若 $\mathfrak K$ 的两个基 $(s)$ 与 $(t)$ 通过变换
\[
(s)=(t)P
\]
相联系，则相应互补基之间有逆变关系
\[
(S)=P^{-1}(T).
\]}

\paragraph{证明。}
设 $e^{(i)}=A_1^{(i)}t_1+\cdots+A_n^{(i)}t_n$ 是单位元的 $n$ 个分量 $e^{(i)}$ 的基表示，其中取定了 $\mathfrak K$ 的一个 $P$-基 $t_1,\ldots,t_n$；该基于是也成为 $\mathfrak K_\Gamma$ 的 $\Gamma$-基。根据第 2、3 小节，$A_1^{(i)},\ldots,A_n^{(i)}$ 属于 $K^{(i)}$，而不同的 $e^{(i)}$ 对应于彼此共轭的系数系 $A$。若 $\alpha_1^{(i)},\ldots,\alpha_n^{(i)}$ 表示 $K^{(i)}$ 中与 $t_1,\ldots,t_n$ 同构对应的基，则还有以下关系：
\begin{equation}
e^{(i)}[t\to\alpha^{(i)}]=e;
\qquad e^{(i)}[t\to\alpha^{(j)}]=0\quad\text{当 }i\ne j. \tag{1}
\end{equation}
因为 $\mathfrak B_{K^{(i)}}^{(i)}$ 包含元素 $t_1-\alpha_1^{(i)},\ldots,t_n-\alpha_n^{(i)}$——它们构成 $\mathfrak B_{K^{(i)}}^{(i)}$ 的一个非独立 $P$-模基——所以在代入 $t=\alpha^{(i)}$ 时该理想消失。于是由
\[
e\equiv e^{(i)}\pmod{\mathfrak B_{K^{(i)}}^{(i)}}
\]
可得 (1) 的第一个关系；第二个关系则由 $i\ne j$ 时 $e^{(i)}$ 属于所有 $\mathfrak B_{K^{(j)}}^{(j)}$ 而得到。按系数展开后，关系式 (1) 可写成矩阵形式：
\begin{equation}
\begin{pmatrix}
\alpha_1^{(1)}&\cdots&\alpha_n^{(1)}\\
\vdots&&\vdots\\
\alpha_1^{(n)}&\cdots&\alpha_n^{(n)}
\end{pmatrix}
\begin{pmatrix}
A_1^{(1)}&\cdots&A_1^{(n)}\\
\vdots&&\vdots\\
A_n^{(1)}&\cdots&A_n^{(n)}
\end{pmatrix}=E_n. \tag{2}
\end{equation}
其中 $E_n$ 表示单位矩阵；由此，$A$ 已由 $\alpha$ 唯一确定。关系式 (2) 表明两个矩阵的行列式都不为零。它还表明，$A_1^{(i)},\ldots,A_n^{(i)}$ 关于 $P$ 线性无关；因为若存在一个线性相关关系，它也会对各共轭系成立，从而降低 $A$ 矩阵的秩，这与 $E_n$ 的秩为 $n$ 相矛盾。因此，每组 $A_1^{(i)},\ldots,A_n^{(i)}$ 都构成 $K^{(i)}$ 的一个基；这些基彼此共轭，所以它们都同 $\mathfrak K$ 中的同一个基 $T_1,\ldots,T_n$ 同构，而根据第 3 小节，$T_\lambda$ 由
\[
T_\lambda=A_\lambda^{(1)}e^{(1)}+\cdots+A_\lambda^{(n)}e^{(n)}
\]
确定。称基 $T_1,\ldots,T_n$ 为 $t_1,\ldots,t_n$ 的互补基，相应地，称 $A_1^{(i)},\ldots,A_n^{(i)}$ 为 $\alpha_1^{(i)},\ldots,\alpha_n^{(i)}$ 的互补基。由于 $A$ 已通过 (2) 唯一确定，而这一关系——通过转到转置矩阵——又从 $T$ 的基、亦即从 $A$ 确定 $\alpha$，所以原来的 $t$ 基就是 $T$ 基的互补基。因此，单位元各分量所给出的对应把所有基都分成互补基对。若 $s,S$ 是另一对互补基，满足
\[
e^{(i)}=B_1^{(i)}s_1+\cdots+B_n^{(i)}s_n
\qquad\text{且}\qquad (s)=(t)P,
\]
则由
\[
A_1t_1+\cdots+A_nt_n=B_1s_1+\cdots+B_ns_n
\]
——在代入 $(s)=(t)P$ 后作为关于 $t$ 的恒等式——可知，$B,A$ 的变换与 $s,t$ 的变换互为逆变。因而由同构可知，$S$ 和 $T$ 也有同样的关系。至此，该定理的各部分全都得证。

\paragraph{注。}
若 $c=C_1t_1+\cdots+C_nt_n$ 是 $\mathfrak K$ 中一个元素关于某个 $P$-基的表示，而 $c=c_1T_1+\cdots+c_nT_n$ 是同一元素关于互补基的表示，则
\[
C_\lambda=\operatorname{Sp}(cT_\lambda),
\qquad c_\lambda=\operatorname{Sp}(ct_\lambda),
\]
其中 $\operatorname{Sp}(x)$ 照例表示各共轭之和：
\[
\operatorname{Sp}(x)=\xi^{(1)}+\xi^{(2)}+\cdots+\xi^{(n)}
\]
对 $\mathfrak K$ 中任意 $x$ 都如此。事实上，
\[
c=\gamma^{(1)}e^{(1)}+\cdots+\gamma^{(n)}e^{(n)}
=\sum_\lambda\gamma^{(1)}A_\lambda^{(1)}t_\lambda+
\cdots+\sum_\lambda\gamma^{(n)}A_\lambda^{(n)}t_\lambda,
\]
由此比较系数便得到 $C_\lambda=\operatorname{Sp}(cT_\lambda)$。第二个关系可同样通过交换 $t$ 与 $T$ 得到。\srcfn{19)}{本小节的讨论表明，不变结构定理实质上比非不变的基定理简单。这里只是为了说明它们同已有结果的联系，才使用后者。事实上，伽罗瓦扩张环的分解以及单位元分量的引入，已经完成了互补基（下一节中的互补模）所完成的一切。下一节中，不同与互补模之间的联系也只是为了归入已有理论才给出。对第 5 小节中拉格朗日插值公式的讨论同样如此。}

\paragraph{5. 以定义方程为基础。拉格朗日插值公式。}
[编者注：本小节在手稿中仍未填写。]
% Source: Paper 43, printed pp. 17--21 (journal scan pp. 706--710)
\subsection*{\S{} 6. 一个阶的不同。}

现在要在整性的意义下加强那些只涉及域及其伽罗瓦扩张环的结构定理：在 (P) 中指定某些子环，并在 \(\mathfrak K\) 和 (K) 中指定关于这些子环为整的元素，也就是各个阶。§ 5 的全部记号继续沿用。\srcfn{20)}{关于元素相对于一个环为整以及整闭的概念，参见脚注 17 所引论文的 § 1。}

\paragraph{1. 基础环。}
设 \(\mathfrak h\) 是 (P) 的一个子环，使 (P) 成为 \(\mathfrak h\) 的分式域；再设 \(\mathfrak h\) 在 (P) 中\emph{整闭}（(P) 中每个相对于 \(\mathfrak h\) 为整的元素都属于 \(\mathfrak h\)）。再令 \(\mathfrak O\) 表示 \(\mathfrak K\) 中的一个 \(\mathfrak h\)-阶，也就是 \(\mathfrak K\) 的一个包含 \(\mathfrak h\) 的子环，其分式域为 \(\mathfrak K\)，并且具有一个有限的、不一定独立的 \(\mathfrak h\)-模基。于是 \(\mathfrak O\) 的所有元素都是（相对于 \(\mathfrak h\)）整的；这个 \(\mathfrak h\)-模基同时成为 \(\mathfrak K\) 的一个——通常相关的——(P)-模基。令 \(\mathfrak o\) 表示 (K) 中与 \(\mathfrak O\) 同构的阶，\(\mathfrak o^{(i)}\) 表示 \(K^{(i)}\) 中的各个共轭阶。\emph{令 \(\mathfrak g\) 为 \(\Gamma\) 中由这 (n) 个共轭阶生成的环；}那么 \(\mathfrak g\) 也是一个阶。事实上，\(\mathfrak g\) 包含 \(\mathfrak h\)。各个 \(\mathfrak o^{(i)}\) 的 \(\mathfrak h\)-基中所有元素的乘积构成 \(\mathfrak g\) 的一个 \(\mathfrak h\)-基，同时构成 \(\Gamma\) 的一个 (P)-基。

\emph{直积 \(\mathfrak O_{\mathfrak o}\) 和 \(\mathfrak O_{\mathfrak g}\) 存在，}这由 § 2, 3 的充分判据得出（其中的环 \(\mathfrak f\) 要用 (P) 替换）。在 \(\mathfrak K_\Gamma\) 中，关于 \(\mathfrak h\) 的直积 \(\mathfrak O_P=\mathfrak K\) 首先存在。因为 \(\mathfrak h\) 包含在交 \([\mathfrak O,P]\) 中；反过来，交中的每个元素作为 \(\mathfrak O\) 的元素都是整的，因而由于 \(\mathfrak h\) 在 (P) 中整闭，它作为 (P) 的元素也属于 \(\mathfrak h\)。取 \(\mathfrak O\) 中任意 (n) 个线性无关元素，它们以 (P) 中元素为系数的线性组合便给出作为直积的 \(\mathfrak K\)。同样有
\[
\mathfrak o_P=K
\qquad\text{以及}\qquad
\mathfrak g_P=\Gamma;
\]
由关于 \(\mathfrak h\) 的直积 \(\mathfrak K_\Gamma=\mathfrak O_P\times\mathfrak o_P\)，按 § 2, 3 的判据可得关于 \(\mathfrak h\) 的 \(\mathfrak O_{\mathfrak o}\) 存在；同样，\(\mathfrak K_\Gamma=\mathfrak O_P\times\mathfrak g_P\) 给出 \(\mathfrak O_{\mathfrak g}\) 的存在。随着 \(\mathfrak O_{\mathfrak o}\) 的存在，各个共轭环 \(\mathfrak O_{\mathfrak o^{(i)}}\) 也存在。称 \(\mathfrak O_{\mathfrak g}\) 为\emph{阶 \(\mathfrak O\) 的伽罗瓦扩张环}。

因此，\(\mathfrak O\) 和 \(\mathfrak o\) 相对于 \(\mathfrak h\) 的不同都有定义；各个共轭不同也同样有定义。令
\[
\mathfrak B=\{\ldots,y-\eta,\ldots\}
\]
表示 \(\mathfrak O_{\mathfrak o}\) 的差分理想，并令
\[
\mathfrak B^{(i)}=\{\ldots,y-\eta^{(i)},\ldots\}
\]
表示 \(\mathfrak O_{\mathfrak o^{(i)}}\) 中的各个共轭差分理想；再令
\[
\mathfrak A^{(i)}=(0):\mathfrak B^{(i)}
\quad\text{——商在 \(\mathfrak O_{\mathfrak o^{(i)}}\) 中取得——}
\]
表示各个共轭\emph{差分商}，特别地，
\[
\mathfrak A=(0):\mathfrak B\quad\text{在 }\mathfrak O_{\mathfrak o}\text{ 中。}
\]
再令
\[
\mathfrak D=\mathfrak A[\eta\to y]
\]
表示 \(\mathfrak O\) 的不同，并令
\[
\mathfrak d^{(i)}=\mathfrak A^{(i)}[y\to\eta^{(i)}]
\]
表示这 (n) 个共轭阶 \(\mathfrak o^{(i)}\) 的不同，特别地，
\[
\mathfrak d=\mathfrak A[y\to\eta]
\]
表示 \(\mathfrak o\) 的不同。\srcfn{21)}{例如，当 \(\mathfrak h\) 是一个数域的主阶，特别是 \(\mathfrak h\) 为全体有理整数所成的系统时，\(\mathfrak h\) 在 (P) 中整闭这一假设成立；因此任意阶的不同和相对不同，包括相对域中的情形，都有定义。再者，\(\mathfrak h\) 可以是 (n) 个不定元的多项式环，这便给出 (n) 个不定元的代数函数域中的不同。}

\paragraph{2. 结构定理。}
与域的伽罗瓦扩张环的结构定理（§ 5, 2）相对应的是：

\paragraph{阶的伽罗瓦扩张环结构定理。}
\emph{相对于阶和域的差分理想与差分商具有如下关系：\(\mathfrak B_K\) 是 \(\mathfrak O_{\mathfrak o}\) 的差分理想 \(\mathfrak B\) 的扩张；同样，\(\mathfrak A_K\) 是 \(\mathfrak O_{\mathfrak o}\) 的差分商 \(\mathfrak A\) 的扩张，而 \(\mathfrak A\) 是 \(\mathfrak A_K\) 的限制。差分商、不同和单位元分量之间具有关系}
\[
\mathfrak A=\mathfrak d e^{(1)},
\qquad
\mathfrak A^{(i)}=\mathfrak d^{(i)}e^{(i)},
\]
\emph{因而}
\[
\mathfrak d=[\mathfrak o,\mathfrak A^{(1)}+\cdots+\mathfrak A^{(n)}].
\]
\emph{所以，\(\mathfrak o\) 的不同 \(\mathfrak d\) 由 (K) 中所有满足 \(xe^{(1)}\in\mathfrak O_{\mathfrak o}\) 的元素 (x) 组成；它不是零理想。}

\(\mathfrak B_K\) 是 \(\mathfrak B\) 的扩张，是因为 \(\mathfrak K=\mathfrak O_P\)，同样 \(K=\mathfrak o_P\)，所以 \(\mathfrak B_K\) 中每个 \(x-\xi\) 都能由 \(\mathfrak B\) 中的各个 \(y-\eta\)，以 \(K\) 中元素为系数作线性表示。由此还可得，\(\mathfrak A_K=(0):\mathfrak B_K\) 也可定义为 \(\mathfrak K_K\) 中的商 \((0):\mathfrak B\)。因此，\(\mathfrak O_{\mathfrak o}\) 中的 \(\mathfrak A=(0):\mathfrak B\) 由 \(\mathfrak O_{\mathfrak o}\) 中属于 \(\mathfrak A_K\) 的全部元素组成，即 \(\mathfrak A=[\mathfrak O_{\mathfrak o},\mathfrak A_K]\)，所以它是 \(\mathfrak A_K\) 的限制。由于 \(\mathfrak A_K=Ke^{(1)}\)，于是 \(\mathfrak A=\mathfrak z e^{(1)}\)，其中 \(\mathfrak z\) 是 \(K\) 中的一个 \(\mathfrak o\)-模，因为 \(\mathfrak A\) 本身也是一个 \(\mathfrak o\)-模。但根据不同的定义，
\[
\mathfrak d=\mathfrak A[x\to\xi]
=(\mathfrak z e^{(1)})[x\to\xi]
=\mathfrak z\bigl(e^{(1)}[x\to\xi]\bigr)=\mathfrak z e=\mathfrak z
\]
（参见 § 5, 4, (1)）。所以 \(\mathfrak z=\mathfrak d\)；差分商 \(\mathfrak A\) 成为 \(\mathfrak o\) 的不同 \(\mathfrak d\) 在 (K) 中的分量，各个共轭情形也相应成立。于是，根据 § 5, 3 的推论，便得到上面关于 \(\mathfrak d\) 的关系。并且
\[
\mathfrak A=\mathfrak d e^{(1)}=[Ke^{(1)},\mathfrak O_{\mathfrak o}]
\]
表明，\(\mathfrak d\) 由 (K) 中所有使 \(xe^{(1)}\) 属于 \(\mathfrak O_{\mathfrak o}\) 的元素 (x) 组成。由此 \(\mathfrak d\) 不是零理想；因为 \(e^{(1)}\) 作为 \(\mathfrak K_K=\mathfrak O_P\times\mathfrak o_P\) 的元素，是 \(\mathfrak O\) 中有限多个元素以 (K) 中系数作成的线性组合，因而是 \(\mathfrak O_{\mathfrak o}\) 中一个元素除以 \(\mathfrak o\) 中一个元素（甚至可以取自 \(\mathfrak h\)）所得的商；这个分母给出 \(\mathfrak d\) 的一个非零元素。于是，\(\mathfrak d\) 关于 (P) 的秩为 (n)，或者说 \(\mathfrak d_P=K\)，从而 \(\mathfrak A_K\) 是 \(\mathfrak A\) 的扩张。

\begin{center}
\textbf{以下为提纲。}
\end{center}

\paragraph{3. 在基独立时同互补模的关系。}
若 \(t_1,\ldots,t_n\) 是 \(\mathfrak O\) 的一个独立 \(\mathfrak h\)-基，那么，只要变换矩阵 (P) 的元素属于 \(\mathfrak h\) 且为幺模，\(\mathfrak K\) 的每个基 \((s)=(t)P\) 也是如此。此时，\((S)=P^{-1}(T)\) 表明，除 (T) 外，(S) 也是同一个 \(\mathfrak h\)-模的基；这个模称为 \(\mathfrak O\) 的互补模。

现在，2 中结构定理的最后一部分说明，\(\mathfrak d\) 等于 (K) 中的商
\[
\mathfrak o:\mathfrak c,
\]
其中 \(\mathfrak c\) 表示互补模。事实上，设
\[
e^{(1)}=A_1t_1+\cdots+A_nt_n,
\]
其中 \(A_i\) 是 \(\mathfrak c\) 的一个基；那么 \(\mathfrak d\) 由所有使 \(\delta e^{(1)}\) 属于 \(\mathfrak O_{\mathfrak o}\) 的元素 \(\delta\) 组成，也就是说，要求 \(\delta A_i\in\mathfrak o\)，所以 \(\mathfrak d=\mathfrak o:\mathfrak c\)。

\paragraph{特殊情形。}
正则阶以 \(e,z,\ldots,z^{n-1}\) 为模基时：
\[
\mathfrak d=\{f'(z)\}.
\]

\paragraph{注。}
由 § 5, 4 的注还可得到 \(\mathfrak O\) 的互补模 \(\mathfrak C\) 的另一个熟知定义。\emph{\(\mathfrak C\) 由 \(\mathfrak K\) 中所有这样的元素 (c) 组成：迹 \(\operatorname{Sp}(ct)\)，也就是任取 \(\mathfrak O\) 中的 (o) 时的迹 \(\operatorname{Sp}(co)\)，都是整的，即属于 \(\mathfrak h\)；}关于 \(\mathfrak o\) 和 (K) 的 \(\mathfrak c\) 也相应如此。

事实上，设 \(c=c_1T_1+\cdots+c_nT_n\)，其中 \(c_i\in\mathfrak h\)，这是 \(\mathfrak C\) 中元素 (c) 关于与 (t) 互补的 \(\mathfrak K\) 的基所作的表示。那么 \(c_i=\operatorname{Sp}(ct_i)\)，故 \(c_i\in\mathfrak h\)；反过来，由 \(\operatorname{Sp}(ct_i)\in\mathfrak h\) 可知 \(c_i\in\mathfrak h\)，因而 \(c\in\mathfrak C\)。

\paragraph{补充。}
\emph{若阶中存在独立的 \(\mathfrak h\)-模基，则差分商 \(\mathfrak A^{(i)}\) 是除 \(\mathfrak B^{(i)}\) 外所有差分理想的交；更确切地说，}
\[
\mathfrak A^{(i)}=
[\mathfrak O_{\mathfrak o^{(i)}},
\mathfrak B_{\mathfrak g}^{(1)},\ldots,
\mathfrak B_{\mathfrak g}^{(i-1)},
\mathfrak B_{\mathfrak g}^{(i+1)},\ldots,
\mathfrak B_{\mathfrak g}^{(n)}].
\]
因为根据 § 2, 2 的加强形式，此时 \(\mathfrak B\) 也是 \(\mathfrak B_K\) 的限制理想；事实上，若 \(e,t_2,\ldots,t_n\) 是 \(\mathfrak O\) 的一个模基，因而也是 \(\mathfrak O_{\mathfrak o}\) 的一个模基，那么
\[
e,t_2-\alpha_2,\ldots,t_n-\alpha_n
\]
也是一个模基；而 \(t_2-\alpha_2,\ldots,t_n-\alpha_n\) 同时是 \(\mathfrak B_K\) 的一个 (K)-模基。因此，§ 2, 2（补充）的假设得到满足。现在，原有
\[
\mathfrak A_\Gamma=[\mathfrak B_\Gamma^{(2)},\ldots,\mathfrak B_\Gamma^{(n)}];
\qquad
\mathfrak A_\mathfrak g=[\mathfrak A_\Gamma,\mathfrak O_\mathfrak g]
=[\ldots,[\mathfrak B_\Gamma^{(i)},\mathfrak O_\mathfrak g],\ldots]
=[\mathfrak B_\mathfrak g^{(2)},\ldots,\mathfrak B_\mathfrak g^{(n)}].
\]
而 \(\mathfrak A=[\mathfrak A_\mathfrak g,\mathfrak O_{\mathfrak o}]\)；所以
\[
\mathfrak A=[\mathfrak O_{\mathfrak o},\mathfrak B_\mathfrak g^{(2)},\ldots,\mathfrak B_\mathfrak g^{(n)}].
\]

\paragraph{4. 在基本方程存在时同其微分商的关系。}
把阶 \(\mathfrak O\) 以 (n) 个不定元 \(u_1,\ldots,u_n\) 扩张，也就是形成直积 \(\mathfrak O_U\)，其中 (U) 表示多项式环 \(\mathfrak h[u_1,\ldots,u_n]\)，故 \(\mathfrak O_U=\mathfrak O[u_1,\ldots,u_n]\)。于是，\(\mathfrak O\) 中每个理想 \(\mathfrak C\) 的扩张 \(\mathfrak C_U\)，由 (u) 的所有系数属于 \(\mathfrak C\) 的多项式组成；反过来，\(\mathfrak C\) 又被刻画为 \(\mathfrak C_U\) 的所有系数所成的系统。

现在假定有一个基本方程；也就是说，当过渡到由 \(\mathfrak h[u]\) 中一个或多个本原多项式 (H(u)) 所确定的商环时，\(e,U,\ldots,U^{n-1}\) 应构成一个 \(\mathfrak h\)-模基。于是，理想 \(\mathfrak C_U\) 过渡为由 \(\mathfrak C_U\) 乘以单位（(H(u)) 及其正负整数次幂）而得到的系统。因此，为了重新得到 \(\mathfrak C\)，可以先通过乘以这些单位回到 \(\mathfrak C_U\)（在 (u) 中为整），再由此回到 \(\mathfrak C\)。特别地，令 \(\mathfrak O\) 的不同 \(\mathfrak D\) 过渡为 \(\mathfrak D_U\)，其中 \(\mathfrak D_U\) 是 \(\mathfrak O_U\) 的不同。但在 \(\mathfrak O_U\) 中——通过用本原的 (H(u)) 取商——存在一个定义方程 (G(U)=0)。所以 (G'(U)) 成为不同的基多项式，也就是说，不同由 (G'(U)) 的系数导出。（还必须更精确地证明，\((G'(U))\) 确实只包含 \(\mathfrak D_U\) 而不包含更多东西；这也许只在 \(\mathfrak o\) 在 (K) 中整闭时成立。因为乘以 (H(u)) 的因子，还可能从 \(\mathfrak D_U\) 中产生更多仍位于 \(\mathfrak D_U\) 内的元素。在整闭的情形下〔五条公理〕，这可由系数理想〔内容〕相乘的定理推出。于是，通过考察基本方程模 (p)，便可建立分歧理论。）

\subsection*{\S{} 7. 主阶模 \(p^t\) 的分歧理论。}

设给定一个数域，\(\mathfrak h\) 为其整数；或者也可以给定一个相对于主阶按某个素理想所得商环的相对域，此时 (p) 表示这个商环的素理想的基元。添入不定元 \(u_1,\ldots,u_n\)；于是，无论主阶还是按 \(p^t\) 取的剩余类环，都具有 \(\mathfrak h\)-模基 \(e,U,\ldots,U^{n-1}\)，因此两种情形中的不同都由 \(G'(u)\) 给出。\emph{主阶的不同模 \(p^t\) 后，过渡为按 \(p^t\) 取的剩余类环的不同}（这对任意阶是否也成立？）。如果
\[
p=\mathfrak p_1^{\varrho_1}\cdots\mathfrak p_r^{\varrho_r},
\]
那么按 \(p^t\) 取的剩余类环 \(\mathfrak R\) 是一个直和，其分量同按 \(\mathfrak p_i^{t\varrho_i}\) 取的剩余类环同构，即
\[
\mathfrak R=\mathfrak R e_1+\cdots+\mathfrak R e_r;
\]
其中，当 \(i\ne k\) 时，\(e_k\) 被 \(\mathfrak p_i^{t\varrho_i}\) 整除。因此，若
\[
\mathfrak d[p^t]=\mathfrak d[p^t]e_1+\cdots+\mathfrak d[p^t]e_r
=\mathfrak d_1+\cdots+\mathfrak d_r
\]
且 (t) 足够大，那么 \(\mathfrak p_i\) 在 \(\mathfrak d_i\) 中出现的幂次，同它在 \(\mathfrak d[p^t]\) 中、因而同它在 \(\mathfrak d\) 中出现的幂次相同。但根据 § 4，\(\mathfrak d_i\) 是 \(\mathfrak R_i\) 的不同；所以只须考察按 \(\mathfrak p^{\varrho t}\) 取的剩余类环的不同。选取 \(\xi\)（《数论报告》定理 29、30），使 \(\xi\) 为整数，并且
\[
\varphi(\xi)\equiv0(p),
\qquad \not\equiv0(p^2),
\]
其中 \(\varphi(x)\) 是模 (p) 的一个 (f) 次素函数。对 \(t\ge2\)，把这个 \(\xi\) 换成 \(\xi^*=\xi e_1\)，其中 \(e_1\) 是剩余类环 \(\mathfrak R\) 模 \(p^t\) 的分解中出现的第一个幂等元；那么，除 \(\varphi(\xi^*)\equiv0(p)\)、\(\not\equiv0(p^2)\) 外，甚至还有 \(\mathfrak p_1=(p,\varphi(\xi^*))\)。于是，关于模 \(p^t\) 的整数的一组基由 \(\xi^\mu\varphi(\xi)^\nu\) 给出，其中 \(\mu=0,1,\ldots,f-1\)，\(\nu=0,1,\ldots,t-1\)。一个独立的 \(\mathfrak h\)-基（模 \(p^t\) 的剩余类）由
\begin{equation}
\xi^{\mu_0}\varphi(\xi)^{\nu_0}
\quad\text{其中}\quad
\mu_0=0,1,\ldots,f-1;
\quad \nu_0=0,1,\ldots,\varrho-1. \tag{1}
\end{equation}
给出。事实上，\((\varphi(\xi),p^{\varrho t})=\mathfrak p\)；所以，由于剩余类环中 \((p)=\mathfrak p^\varrho\)，从主理想过渡到元素便有 \(\varphi(\xi)=pM(\xi)\)，其中 \(M(\xi)\) 是剩余类环中的单位，也就是说 \(M(\xi)\not\equiv0(p)\)。因此（由于 \(p^t=0\)），可以逐次用较低次幂表示 \(\xi^{\varrho t}\)，先在 \(M(\xi)\) 中，再在剩余类环中表示；所以 (1) 确实是一组基。关于这组基的定义方程为
\[
F(x)=\varphi(x)^\varrho+pM(x),
\]
其中 (M(x)) 模 (p) 不被 \(\varphi(x)\) 整除（Ö. Ore, Math. Ann. 96, S. 348），因而
\[
M(\xi)\not\equiv0(p).
\]
但这组基也是独立的；因为由
\[
a_0(\xi)+a_1(\xi)\varphi(\xi)+\cdots+a_{\varrho-1}(\xi)\varphi(\xi)^{\varrho-1}
\equiv0(p^{t\varrho})
\]
先可推出模 (p) 同余于零；从而 \(a_0(\xi)\equiv0(p)\)，于是
\[
a_0(x)\equiv0(p);
\]
约去 (p) 后逐次得到
\[
a_0(x)\equiv0(p),\ldots,a_{\varrho-1}(x)\equiv0(p),
\]
再约去 (p) 后得到
\[
a_0(x)\equiv0(p^2),\ldots,a_{\varrho-1}(x)\equiv0(p^2),
\]
如此继续，直至
\[
a_0(x)\equiv0(p^t),\ldots,a_{\varrho-1}(x)\equiv0(p^t).
\]
于是，剩余类环的不同由多项式微分商 \(F'(\xi)\) 给出——全部问题都归结为 Ore, S. 345, § 3。由 Ore 定理 10（参见 Dedekind--Hensel）可知，对每个 \(\mathfrak p\)，取 \(p^t=p^n\) 已经足够；并且还可由 Ore 得到辅助数。

通过过渡到商环，所有这些同时也适用于相对不同；从而 Ore, Math. Ann. 97, S. 594, § 4 直接由 Ore, Math. Ann. 96 得出。只须补充关于如下可表示性的断言：上域的不同等于下域的不同与相对不同的乘积（参见本文引言）。

\begin{center}
1949 年 10 月 25 日收稿。
\end{center}

\end{document}
